\exercisesection{Spectrum of Endomorphisms of Banach Spaces}

In the exercises below, all Banach spaces are over C.

\begin{enumerate}
\def\labelenumi{\arabic{enumi})}
\tightlist
\item
  Let \(\mathrm{E} = \mathcal{C}([0,1])\) and let \(u\) be the endomorphism of \(\mathrm{E}\) defined by
\end{enumerate}

\[
u (f) (t) = \int_ {0} ^ {t} f (t) d t.
\]

\begin{enumerate}
\def\labelenumi{\alph{enumi})}
\item
  The spectral radius of \(u\) is zero. (Estimate the norm of \(u^n\) .)
\item
  The spectrum of \(u\) consists only of 0, and, for every \(n \geqslant 1\), the kernel of \(u^{n}\) consists only of 0.
\end{enumerate}

(In particular, 0 is an isolated point of the spectrum of \(u\), but the closure of the union of the kernels of \(u^n\), for \(n \geqslant 1\), is not equal to the spectral subspace \(\mathrm{E}_0(x)\).)

\begin{enumerate}
\def\labelenumi{\arabic{enumi})}
\setcounter{enumi}{1}
\item
  There exists a Banach space E and an endomorphism u of E such that 0 is an eigenvalue of u but 0 is not isolated in \(\mathrm{Sp}(u)\).
\item
  Give an example of a Hilbert space E and an endomorphism \(u \in \mathcal{L}(\mathrm{E})\) such that 0 is an eigenvalue of u but not of \(u^{*}\).
\item
  Let \(E = \ell^{2}(\mathbf{Z})\) and let \(u\) be the endomorphism such that \(u((\lambda_{n})) = (\lambda_{n+1})\).
\end{enumerate}

\begin{enumerate}
\def\labelenumi{\alph{enumi})}
\item
  Show that \(u\) is unitary.
\item
  Compute the spectrum of \(u\).
\item
  The finite spectrum of \(u\) is empty.
\end{enumerate}

\begin{enumerate}
\def\labelenumi{\arabic{enumi})}
\setcounter{enumi}{4}
\tightlist
\item
  Let \(s\) be the endomorphism of \(\ell^2 (\mathbf{N})\) such that
\end{enumerate}

\[
s (a _ {0}, a _ {1}, \dots , a _ {n}, \dots) = (a _ {1}, \dots , a _ {n}, \dots).
\]

Show that the numerical range of \(s\) is the open unit disk \(\Delta\) in C, and that the spectrum of \(s\) is the closed unit disk in C. Determine the eigenvalues of \(s\). Do the same for the adjoint \(s^{*}\) of \(s\).

\begin{enumerate}
\def\labelenumi{\arabic{enumi})}
\setcounter{enumi}{5}
\item
  Let \(u \in \mathcal{L}(\mathbf{C}^n)\) be a trace-zero endomorphism. There exists an orthonormal basis B of \(\mathbf{C}^n\) such that all diagonal entries of the matrix representing \(u\) in the basis B are zero. (Prove that 0 belongs to the numerical range of \(u\).)
\item
  Let \(u\) be an endomorphism of a complex Hilbert space E such that the spectrum of \(u\) is contained in the open unit disk in C. Let F be the Hilbert space \(\ell_{\mathrm{E}}^{2}(\mathbf{N})\) of square-summable sequences of elements of E (EVT, V, p.~18, example), and let \(s_{\mathrm{F}}\) be the endomorphism of F such that
\end{enumerate}

\[
s _ {\mathrm{F}} (a _ {0}, a _ {1}, \ldots , a _ {n}, \ldots) = (a _ {1}, \ldots , a _ {n} \ldots).
\]

There exists a closed subspace H of F and a continuous bijective linear map \(v: E \rightarrow H\) such that \(u = v^{-1}s_{F}v\).

\begin{enumerate}
\def\labelenumi{\arabic{enumi})}
\setcounter{enumi}{7}
\tightlist
\item
  Let \(u\) be an endomorphism of a complex Hilbert space E. Let \(S_u\) be the set of endomorphisms of E of the form \(wuw^{-1}\), where \(w\) ranges over the set of invertible endomorphisms of E.
\end{enumerate}

\begin{enumerate}
\def\labelenumi{\alph{enumi})}
\tightlist
\item
  We have
\end{enumerate}

\[
\varrho (u) = \inf _ {v \in \mathrm{S} _ {u}} \| v \|.
\]

(One may assume that \(\varrho(u) \leqslant 1\); apply the preceding exercise to the endomorphisms \(u_{\varepsilon} = (1 + \varepsilon)^{-1} u\) for \(\varepsilon > 0\).)

\begin{enumerate}
\def\labelenumi{\alph{enumi})}
\setcounter{enumi}{1}
\tightlist
\item
  The convex hull of the spectrum of \(u\) is equal to the intersection of the subsets \(\iota(v)\), for \(v \in S_u\) ("Hildebrandt's theorem").
\end{enumerate}

¶ 9) Let E be a complex Banach space.

\begin{enumerate}
\def\labelenumi{\alph{enumi})}
\item
  Let \(u \in \mathcal{L}(\mathrm{E})\). Consider the endomorphisms \(\gamma_{u}\) and \(\delta_{u}\) of \(\mathcal{L}(\mathrm{E})\) defined by \(\gamma_{u}(v) = u \circ v\) and \(\delta_{u}(v) = v \circ u\). We have \(\operatorname{Sp}(u) = \operatorname{Sp}(\gamma_{u}) = \operatorname{Sp}(\delta_{u})\), where the last two spectra are relative to the algebra \(\mathcal{L}(\mathcal{L}(\mathrm{E}))\).
\item
  For every function \(f\) holomorphic in a neighborhood of the spectrum of \(u\), one has \(\boldsymbol{\gamma}_{f(u)} = f(\boldsymbol{\gamma}_{u})\) and \(\boldsymbol{\delta}_{f(u)} = f(\boldsymbol{\delta}_{u})\).
\item
  Let \(u_{1}\) and \(u_{2}\) be in \(\mathcal{L}(\mathrm{E})\). Let \(\varrho\) be the endomorphism \(v \mapsto u_{1} \circ v \circ u_{2}\) of \(\mathcal{L}(\mathrm{E})\). Its spectrum is equal to \(\operatorname{Sp}(u_{1})\operatorname{Sp}(u_{2}) \subset \mathbf{C}\).
\end{enumerate}

\(d)\) Let \(u_{1}\) and \(u_{2}\) be in \(\mathcal{L}(\mathrm{E})\). Let \(\varrho\) be the endomorphism \(v \mapsto u_{1} \circ v \circ u_{2}\) of the space \(\mathcal{L}_{2}(\mathrm{E})\) of Hilbert--Schmidt maps from E to E. Its spectrum is equal to \(\mathrm{Sp}(u_{1})\mathrm{Sp}(u_{2}) \subset \mathbf{C}\). (Apply a similar method.)

\begin{enumerate}
\def\labelenumi{\arabic{enumi})}
\setcounter{enumi}{9}
\item
  Let E be a complex Hilbert space. Let \(x_1\) and \(x_2\) be endomorphisms of E. Show that the spectrum of the endomorphism \(x_1 \otimes x_2\) of E \(\widehat{\otimes}_2\) E is equal to \(\mathrm{Sp}(x_1) \mathrm{Sp}(x_2)\). (Identify \(\overline{\mathrm{E}} \widehat{\otimes}_2\) E with the space of Hilbert--Schmidt operators from E to E, cf.~EVT 5, p.~52, th. 1, b), and apply the preceding exercise.)
\item
  Let \(u\) be an endomorphism of a complex Hilbert space \(E\) and let \((j, |u|)\) be its polar decomposition.
\end{enumerate}

\begin{enumerate}
\def\labelenumi{\alph{enumi})}
\item
  For \(j\) to commute with \(|u|\), it is necessary and sufficient that \(u\) commute with \(u^{*}u\).
\item
  Prove that one then has \(\|u^{*}(x)\|\leqslant\|u(x)\|\) for all \(x\in E\), and that \(\varrho(u)=\|u\|\).
\item
  Let \(E = \ell^{2}(\mathbf{N})\) and let \(u\) be the endomorphism of E given by
\end{enumerate}

\[
u ((a _ {0}, a _ {1}, \dots)) = (0, a _ {0}, a _ {1}, \dots).
\]

It is not normal. Compute the polar decomposition \((j, |u|)\) of u; \(j\) commutes with \(|u|\).

\begin{enumerate}
\def\labelenumi{\arabic{enumi})}
\setcounter{enumi}{11}
\tightlist
\item
  Let \(u\) be an endomorphism of a Banach space E. We denote by \(\sigma_{a}(u)\) the set of \(\lambda\in\mathbf{C}\) such that, for every \(\alpha\in\mathbf{R}_{+}^{*}\), there exists \(x\in\mathrm{E}\) such that \(\|\lambda x-u(x)\|<\alpha\|x\|\).
\end{enumerate}

\begin{enumerate}
\def\labelenumi{\alph{enumi})}
\item
  Prove that \(\sigma_{a}(u)\) is a closed subset of \(\sigma(u)\).
\item
  Prove that \(\sigma_{a}(u)\) contains the boundary of \(\sigma(u)\).
\end{enumerate}

\begin{enumerate}
\def\labelenumi{\arabic{enumi})}
\setcounter{enumi}{12}
\tightlist
\item
  Let E be a complex Hilbert space. Let N be a map from \(\mathcal{L}(\mathrm{E})\) into \(\mathfrak{P}(\mathbf{C})\) such that
\end{enumerate}

\begin{enumerate}
\def\labelenumi{\alph{enumi})}
\item
  \(N(u)\) is compact for every \(u \in \mathcal{L}(\mathrm{E})\);
\item
  \(\mathrm{N}(\alpha u + \beta 1_{\mathrm{E}}) = \alpha \mathrm{N}(u) + \beta\) for all \(u \in \mathcal{L}(\mathrm{E})\) and every \((\alpha, \beta) \in \mathbf{C}^{2}\);
\item
  \(\mathrm{N}(u)\subset\{z\in\mathbf{C}\mid\mathcal{R}(z)\geqslant0\}\) if and only if \(u+u^{*}\) is positive.
\end{enumerate}

  Then \(\mathrm{N}(u)=\overline{\iota(u)}\) for all \(u\in\mathcal{L}(\mathrm{E})\). (Show that \(u\mapsto\overline{\iota(u)}\) has the indicated properties; let \(N_{1}\) and \(N_{2}\) be maps satisfying them; assuming that \(\mathrm{N}_{1}(u)\) is not contained in \(\mathrm{N}_{2}(u)\), find \(\alpha\) and \(\beta\) in C such that \(\mathrm{N}_{2}(\alpha u+\beta)\) is contained in the half-plane \(\mathcal{R}(z)\geqslant0\) but \(\mathrm{N}_{1}(\alpha u+\beta)\) is not, and conclude a contradiction using c).)

\begin{enumerate}
\def\labelenumi{\arabic{enumi})}
\setcounter{enumi}{13}
\item
  Give an example of a complex Hilbert space and of \(u \in \mathcal{L}(\mathrm{E})\) such that \(\iota(u)\) is not closed.
\item
  The set X of nonempty compact subsets of C is equipped with the distance defined in TG, IX, p.~91, exerc. 6. Let E be a complex Hilbert space.
\end{enumerate}

\begin{enumerate}
\def\labelenumi{\alph{enumi})}
\item
  The map \(u \mapsto \overline{\iota(u)}\) from \(\mathcal{L}(\mathrm{E})\) to X is continuous. If E is infinite-dimensional, it is not continuous when \(\mathcal{L}(\mathrm{E})\) is equipped with the topology of simple convergence.
\item
  The map \(u \mapsto \operatorname{Sp}(u)\) from \(\mathcal{L}(\mathrm{E})\) to X is not continuous.
\end{enumerate}

\begin{enumerate}
\def\labelenumi{\arabic{enumi})}
\setcounter{enumi}{15}
\tightlist
\item
  Let X be a locally compact space, \(\mu\) a positive measure of total mass 1 on X, and \(f: X \to X\) a map such that \(f(\mu) = \mu\).
\end{enumerate}

\begin{enumerate}
\def\labelenumi{\alph{enumi})}
\item
  The map \(\varphi \mapsto \varphi \circ f\) induces a unitary endomorphism \(u_{f}\) of \(\mathrm{L}^2 (\mathrm{X},\mu)\).
\item
  The real number 1 is an eigenvalue of \(u_f\) of multiplicity 1 if and only if, for every measurable subset A of X such that \(f(A) = A\), one has \(\mu(A) \in \{0,1\}\). One then says that \(f\) is \(\mu\)-ergodic.
\item
  Let \(\mathrm{E}\subset\mathrm{L}^{2}(\mathrm{X},\mu)\) be the eigenspace of \(u_{f}\) for the eigenvalue 1. The sequence \((v_{\mathrm{N}})_{\mathrm{N}\in\mathbf{N}}\) defined by
\end{enumerate}

\[
v _ {\mathrm{N}} = \frac {1}{\mathrm{N}} \sum_ {n = 0} ^ {\mathrm{N-1}} u _ {f} ^ {n}
\]

converges to the orthogonal projection onto E in the space \(\mathcal{L}(\mathrm{L}^{2}(\mathrm{X},\mu))\) equipped with the topology of simple convergence ("von Neumann's ergodic theorem"). (Note that the orthogonal complement of E is the closure of the image of \(u_{f}-1_{\mathrm{E}}\); then verify that \(v_{\mathrm{N}}(\varphi)\to0\) when \(\varphi\in\mathrm{E}^{\circ}\).)

\(d)\) One defines an endomorphism \(\widetilde{u}_f\) of \(\mathrm{L}^1 (\mathrm{X},\mu)\) as in part \(a)\), and we set

\[
\widetilde {v} _ {\mathrm{N}} = \frac {1}{\mathrm{N}} \sum_ {n = 0} ^ {\mathrm{N-1}} \widetilde {u} _ {f} ^ {n}.
\]

For every \(\varphi\in\mathrm{L}^{1}(\mathrm{X},\mu)\) the sequence \((\widetilde{v}_{\mathrm{N}}(\varphi))_{\mathrm{N}\in\mathbf{N}}\) converges in \(\mathrm{L}^{1}(\mathrm{X},\mu)\) to an element \(\widetilde{\varphi}\) of the eigenspace of \(\widetilde{u}_{f}\) for the eigenvalue 1. (First consider the case where \(\varphi\) is bounded, and apply the previous question.)

\begin{enumerate}
\def\labelenumi{\alph{enumi})}
\setcounter{enumi}{4}
\tightlist
\item
  Let I be a finite set and X = I \(^{Z}\) endowed with the product measure of the measures \(\mu_n\) on I such that \(\mu_n(i) = 1/\text{Card(I)}\) for all \(i \in \text{I}\). Set \(f((i_n)) = (i_{n+1})\) for all \((i_n) \in \text{X}\). Then \(f(\mu) = \mu\) and \(f\) is \(\mu\)-ergodic (“Bernoulli shift”).
\end{enumerate}

\(f)\) Let \(\mathrm{X} = \mathbf{R} / \mathbf{Z}\) and \(\mu\) the normalized Haar measure on \(\mathrm{X}\). Let \(a \in \mathrm{X}\) and set \(f(x) = x + a\). Then \(f(\mu) = \mu\) and \(f\) is \(\mu\)-ergodic if and only if \(a\) is the class of an irrational number.

\exercisesection{Algebras of Continuous Functions on a Compact Space}

All Banach spaces and all Banach algebras below are complex.

\begin{enumerate}
\def\labelenumi{\arabic{enumi})}
\tightlist
\item
  Let X be a topological space, let \(\mathcal{C}_{b}(\mathrm{X})\) be the stellar algebra of bounded continuous complex functions on X, and let \(\widehat{\mathrm{X}} = \mathrm{X}(\mathcal{C}_{b}(\mathrm{X}))\).
\end{enumerate}

\begin{enumerate}
\def\labelenumi{\alph{enumi})}
\tightlist
\item
  Let A be a unital Banach subalgebra of \(\mathcal{C}_{b}(\mathrm{X})\) and let \(R_{A}\) be the equivalence relation on X such that x is equivalent to \(x'\) if \(f(x) = f(x')\) for every function \(f \in A\). The \(R_{A}\)-saturation S of a compact subset K of X is closed. (Let \(x \in \overline{S}\); for all \(f \in A\) and every \(\varepsilon > 0\), let \(V_{f,\varepsilon}\) be the set of \(x' \in X\) such that \(|f(x) - f(x')| \leqslant \varepsilon\); if \(f_{1}, \ldots, f_{n} \in A\) and \(\varepsilon_{1}, \ldots, \varepsilon_{n} > 0\), one has
\end{enumerate}

\[
\mathrm{V} _ {f _ {1}, \varepsilon_ {1}} \cap \dots \cap \mathrm{V} _ {f _ {n}, \varepsilon_ {n}} \cap \mathrm{K} \neq \varnothing ,
\]

hence

\[
\bigcap_{\substack{f\in \mathrm{A}\\ \varepsilon >0}}\mathrm{V}_{f,\varepsilon}\cap \mathrm{K}\neq \varnothing
\]

whence \(x \in S\).)

\begin{enumerate}
\def\labelenumi{\alph{enumi})}
\setcounter{enumi}{1}
\item
  Assume X is normal. Let R be a closed equivalence relation on X and let \(A = A_{R}\) be the subalgebra of \(\mathcal{C}_{b}(X)\) formed by the functions constant on the classes modulo R. Show that \(R = R_{A}\) (Use TG, IX, p.~105, exerc. 19, a).
\item
  Assume X is compact. \(R \mapsto A_{R}\) is a bijection from the set of separated equivalence relations on X onto the set of stellar subalgebras of \(\mathcal{C}(X)\) (Let A be a stellar subalgebra of \(\mathcal{C}(X)\). Show that \(R_{A}\) is separated, and apply the Weierstrass–Stone theorem to the algebra of continuous functions on \(X/R_{A}\) obtained from A by passing to the quotient.)
\end{enumerate}

\begin{enumerate}
\def\labelenumi{\arabic{enumi})}
\setcounter{enumi}{1}
\tightlist
\item
  Let X be a compact topological space, R a separated equivalence relation on X, and B a closed subalgebra of \(\mathcal{C}(\mathrm{X})\) containing \(\mathrm{A}_{\mathrm{R}}\) (notation of Exercise 1). If \(f \in \mathcal{C}(\mathrm{X})\) coincides on each R-class \(c\) with an element \(g_{c}\) of B, then \(f \in \mathrm{B}\). (Let \(\varepsilon > 0\). There exists a saturated open neighborhood \(\mathrm{V}_{c}\) of \(c\) such that \(|f - g_{c}| \leqslant \varepsilon\) on \(\mathrm{V}_{c}\). Let \(\mathrm{V}_{c_{1}}, \ldots, \mathrm{V}_{c_{n}}\) be a cover of X. Let \((u_{1}, \ldots, u_{n})\) be a partition of unity subordinate to \((\mathrm{V}_{c_{1}}, \ldots, \mathrm{V}_{c_{n}})\) and consisting of functions of \(\mathrm{A}_{\mathrm{R}}\). Then
\end{enumerate}

\[
g = u _ {1} g _ {c _ {1}} + \dots + u _ {n} g _ {c _ {n}} \in \mathrm{B},
\]

and \(|f - g| \leqslant \varepsilon\).)

\begin{enumerate}
\def\labelenumi{\arabic{enumi})}
\setcounter{enumi}{2}
\tightlist
\item
  Let X be a compact topological space and, for every \(x \in X\), let \(A_{x}\) be a commutative unital Banach algebra, with \(e_{x}\) denoting the identity element. We assume \(\|e_{x}\| = 1\). Let B be the product Banach algebra of the \(A_{x}\) for \(x \in X\). Let C be a Banach subalgebra of B having the following properties:
\end{enumerate}

\begin{enumerate}
\def\labelenumi{(\roman{enumi})}
\item
  \(e = (e_x)_{x \in \mathbf{X}} \in \mathbf{C};\)
\item
  If \(a = (a_x) \in \mathrm{C}\) and \(f \in \mathcal{C}(\mathrm{X})\) , the function \(fa: x \mapsto f(x)a_x\) is an element of \(\mathrm{C}\) ;
\item
  For every \(a = (a_x) \in \mathrm{C}\), the function \(x \mapsto \|a_x\|\) is upper semicontinuous.
\end{enumerate}

Let \(\widehat{\mathbf{X}}\) be the disjoint union of the \(\mathsf{X}(\mathsf{A}_x)\). For every \(x_0 \in \mathbf{X}\) and \(\chi \in \mathsf{X}(\mathsf{A}_{x_0})\), let \(\varphi(\chi)\) be the character of C defined by \(\varphi(\chi)((a_x)_{x \in \mathbf{X}}) = \chi(a_{x_0})\). Prove that \(\varphi\) defines a bijection from \(\widehat{\mathbf{X}}\) onto \(\mathsf{X}(\mathbf{C})\). (Let \(\xi \in \mathsf{X}(\mathbf{C})\). Since \(\mathcal{C}(\mathbf{X})\) can be identified with a closed subalgebra of C, \(\xi\) defines a character of \(\mathcal{C}(\mathbf{X})\), hence there exists \(x_0 \in \mathbf{X}\). Let \(a = (a_x) \in \mathbf{C}\) and \(b = (b_x) \in \mathbf{C}\), with \(a_{x_0} = b_{x_0}\); let \(\varepsilon > 0\); we have \(\|a_x - b_x\| < \varepsilon\) on a neighborhood U of \(x_0\); let \(f \in \mathcal{C}(\mathbf{X})\) be such that \(0 \leqslant f \leqslant 1\), equal to 1 at \(x_0\) and to 0 outside U; we have \(\|fa - fb\| \leqslant \varepsilon\), whence \(|\xi(fa - fb)| \leqslant \varepsilon\), whence \(|\xi(a) - \xi(b)| \leqslant \varepsilon\), whence \(\xi(a) = \xi(b)\). Deduce that there exists a \(\chi \in \mathsf{X}(\mathsf{A}_{x_0})\) such that \(\xi = \varphi(\chi)\).

\begin{enumerate}
\def\labelenumi{\arabic{enumi})}
\setcounter{enumi}{3}
\tightlist
\item
  Let X be a compact topological space such that the Banach algebra \(\mathcal{C}(\mathrm{X})\) is generated by a sequence \((j_n)\) of idempotents. The function
\end{enumerate}

\[
\omega \mapsto \sum_ {n = 1} ^ {\infty} \frac {1}{3 ^ {n}} (2 j _ {n} (\omega) - 1)
\]

separates the points of X. It generates the Banach algebra \(\mathcal{C}(\mathrm{X})\).

\begin{enumerate}
\def\labelenumi{\arabic{enumi})}
\setcounter{enumi}{4}
\tightlist
\item
  Let A be the Banach algebra of functions from {[}0,1{]} to C admitting continuous derivatives on {[}0,1{]} up to order n (example 4 of I, p.~18). Let \(\omega \in [0,1] = \mathsf{X}(\mathsf{A})\).
\end{enumerate}

\begin{enumerate}
\def\labelenumi{\alph{enumi})}
\item
  The smallest closed ideal I of A such that \(V(I) = \{\omega\}\) is the set of \(f \in A\) such that \(f, f', \ldots, f^{(n)}\) vanish at \(\omega\). (Use prop. 5 of I, p.~94.)
\item
  If \(g \in A\), the norm of its canonical image in A/I is equal to
\end{enumerate}

\[
\sum_ {k = 0} ^ {n} \frac {| g ^ {(k)} (\omega) |}{k !}.
\]

\begin{enumerate}
\def\labelenumi{\alph{enumi})}
\setcounter{enumi}{2}
\tightlist
\item
  The algebra A/I is isomorphic to \(\mathbf{C}[\mathrm{X}]/(\mathrm{X}^{n+1})\); the only maximal ideal of A containing I is the set of \(f \in A\) such that \(f(\omega) = 0\).
\end{enumerate}

\begin{enumerate}
\def\labelenumi{\arabic{enumi})}
\setcounter{enumi}{5}
\tightlist
\item
  Let \(\Delta\) be the disk \(|\zeta| \leqslant 1\) in \(\mathbf{C}\) and A the Banach algebra of continuous complex-valued functions on \(\Delta\) that are analytic in the interior of \(\Delta\) (example 9 of I, p.~20).
\end{enumerate}

\begin{enumerate}
\def\labelenumi{\alph{enumi})}
\tightlist
\item
  If \(x \in A\), then \(x\) is the limit in A of the functions
\end{enumerate}

\[
\zeta \mapsto x _ {n} (\zeta) = x \left(\frac {n \zeta}{n + 1}\right).
\]

The identity mapping of \(\Delta\) is a generator of A.

\begin{enumerate}
\def\labelenumi{\alph{enumi})}
\setcounter{enumi}{1}
\item
  The space X(A) is canonically identified with \(\Delta\) (Use assertion g) of prop. 1 of I, p.~142.)
\item
  Let S be a closed subset of \(\Delta\), identified with a closed subset of \(\mathsf{X}(\mathsf{A})\). If S has a non-isolated point \(\zeta\) such that \(|\zeta| < 1\), the closure of S for the Jacobson topology is \(\Delta\).
\end{enumerate}

\(d)\) For \(x \in \mathrm{A}\), set \(x^{*}(\zeta) = x(\overline{\zeta})\). Then \(x \mapsto x^{*}\) is an isometric involution of A. The Hermitian characters of A are given by the real elements of \(\Delta\).

\begin{enumerate}
\def\labelenumi{\alph{enumi})}
\setcounter{enumi}{4}
\tightlist
\item
  The map \(x \mapsto x|U\) is an isomorphism from A onto \(\mathrm{P}(\mathbf{U})\). Every function in \(\mathcal{C}_{\mathbf{R}}(\mathbf{U})\) is the uniform limit of real parts of functions in \(\mathrm{P}(\mathbf{U})\).
\end{enumerate}

\begin{enumerate}
\def\labelenumi{\arabic{enumi})}
\setcounter{enumi}{6}
\tightlist
\item
  Let \(A_{1}\) (resp. \(A_{2}\)) be the unital Banach algebra considered in Exercise 5 (resp. 6). Let \(A = A_{1} \times A_{2}\). Show that A is semisimple, but that the image of the Gelfand transform of A is neither closed nor dense in \(\mathcal{C}(\mathsf{X}(\mathsf{A}))\).
\end{enumerate}

¶ 8) Let X be a compact topological space and B a unital Banach subalgebra of \(\mathcal{C}(\mathrm{X})\) (with the induced norm), separating the points of X. One identifies X with a closed subset of \(\mathrm{X(B)} = \mathrm{X}'\). Let \(\mathrm{B}^*\) be the set of invertible elements of B. One says that B is logmodular if the set of functions \(\log(|f|)\), for \(f \in \mathrm{B}^*\), is dense in \(\mathcal{C}_{\mathbf{R}}(\mathrm{X})\). (This is the case, in particular, if the set of functions \(\mathcal{R}(f)\), for \(f \in \mathrm{B}\), is dense in \(\mathcal{C}_{\mathbf{R}}(\mathrm{X})\), because \(\log(|e^f|) = \mathcal{R}(f)\).) In what follows, one assumes that B is logmodular.

\begin{enumerate}
\def\labelenumi{\alph{enumi})}
\tightlist
\item
  For every \(\chi \in \mathrm{X}'\) , there exists a unique measure \(\mu_{\chi} \geqslant 0\) on X such that
\end{enumerate}

\[
\log (| \chi (f) |) = \int_ {\mathrm{X}} \log (| f (\omega) |) d \mu_ {\chi} (\omega)
\]

for all \(f \in \mathrm{B}^*\), hence such that \(\chi(f) = \int_{\mathrm{X}} f(\omega) d\mu_{\chi}(\omega)\) for all \(f \in \mathrm{B}\). \(b)\) One has

\[
\log | \chi (f) | \leqslant \int_ {\mathrm{X}} \log | f (\omega) | d \mu_ {\chi} (\omega)
\]

for all \(f \in \mathrm{B}\) and the condition \(\chi(f) = \int_{\mathrm{X}} f(\omega) d\mu_{\chi}(\omega)\) for all \(f \in \mathrm{B}\) defines \(\mu_{\chi}\) uniquely.

\begin{enumerate}
\def\labelenumi{\alph{enumi})}
\setcounter{enumi}{2}
\tightlist
\item
  Let \(\chi \in X'\) and \(\mu\) be a positive measure on \(X\). Write \(\mu = \mu_1 + \mu_2\), where \(\mu_1\) is absolutely continuous with respect to \(\mu_{\chi}\), and where \(\mu_2\) and \(\mu_{\chi}\) are mutually singular. Let I be the kernel of \(\chi\). Then
\end{enumerate}

\[
\inf _ {f \in \mathrm{I}} \int_ {\mathrm{X}} | 1 - f | ^ {2} d \mu = \inf _ {f \in \mathrm{I}} \int_ {\mathrm{X}} | 1 - f | ^ {2} d \mu_ {1}.
\]

\(d)\) Let \(\chi\in\mathrm{X}^{\prime}\) and let \(h\) be a nonnegative function in \(\mathcal{L}^{1}(\mu_{\chi})\). Let I be the kernel of \(\chi\). Then

\[
\inf _ {f \in \mathrm{I}} \int_ {\mathrm{X}} | 1 - f | ^ {2} h d \mu_ {\chi} = \exp \Bigl (\int_ {\mathrm{X}} \log (h) d \mu_ {\chi} \Bigr).
\]

\begin{enumerate}
\def\labelenumi{\alph{enumi})}
\setcounter{enumi}{4}
\tightlist
\item
  For every positive measure \(\mu\) on X, and every \(p \in [1, +\infty[\), let \(\mathrm{H}^p(\mu)\) be the closure of the canonical image of B in \(\mathrm{L}^p(\mu)\). Let \(\chi \in \mathrm{X}'\) and let I be the kernel of \(\chi\). Then \(\mathrm{L}^2(\mu_\chi)\) is the Hilbert sum of \(\mathrm{H}^2(\mu_\chi)\) and of the closure in \(\mathrm{L}^2(\mu_\chi)\) of the canonical image of \(\bar{\mathrm{I}}\) (the set of conjugates of functions belonging to I).
\end{enumerate}

\(f)\) The space \(\mathrm{H}^1 (\mu_\chi)\) is the set of classes \(\widetilde{h}\) of \(h\in \mathcal{L}^1 (\mu_\chi)\) such that \(\int_{\mathrm{X}}fh  d\mu_{\chi} = 0\) for all \(f\in \mathrm{I}\).

\(g)\) Let \(h\) be a nonnegative function in \(\mathcal{L}^1 (\mu_\chi)\). Then we have \(\widetilde{h} = |\widetilde{f}|\) with \(\widetilde{f}\in \mathrm{H}^1 (\mu_\chi)\) and \(\int_{\mathrm{X}}f  d\mu_{\chi}\neq 0\) if and only if \(\log (h)\in \mathcal{L}^1 (\mu_\chi)\).

¶ 9) Let X be a compact topological space and B a unital Banach subalgebra of \(\mathcal{C}(\mathrm{X})\) (with the induced norm). Every element \(x\) of X defines a character \(\chi_x\) of B. If \(x\) and \(x'\) are in X, we write \(x \sim x'\) when \(\| \chi_x - \chi_{x'} \| \neq 2\). a) Let \(x\) and \(x'\) be in X. If there exists a sequence \((f_n)\) of elements of B such that \(\| f_n \| \leqslant 1\), \(|f_n(x)|\) tends to 1, and \(\liminf_{n \to \infty} |f_n(x) - f_n(x')| > 0\), then \(x\not\sim x'\). (One may assume that \(f_n(x)\) tends to 1. Consider \(g_n = (f_n - \lambda_n)(1 - \lambda_n f_n)^{-1}\), where \((\lambda_n)\) is a sequence of numbers in \([0,1[\) tending toward 1. We have \(\| g_n \| \leqslant 1\). If the sequence \((\lambda_n)\) is well chosen, \(g_n(x)\) tends to 1 and \(g_n(x')\) tends to -1.)

\begin{enumerate}
\def\labelenumi{\alph{enumi})}
\setcounter{enumi}{1}
\item
  Let R be the set of real parts of the elements of B. If \(x \sim x'\), there exists \(c \in ]0,1[\) such that \(f(x) \geqslant cf(x')\) and \(f(x') \geqslant cf(x)\) for every \(f \geqslant 0\) in R. (It suffices to obtain the first condition; if it is not true, there exists a sequence \((f_n)\) in R with \(f_n \geqslant 0\), \(f_n(x') = 1\), and \(f_n(x)\) tending to 0; let \(g_n \in B\) with \(\mathcal{R}(g_n) = f_n\). Then \(e^{-g_n} \in \mathrm{B}\), \(\| e^{-g_n} \| \leqslant 1\), \(|e^{-g_n(x')}| = 1 / e\), and \(|e^{-g_n(x)}|\) tends to 1, which contradicts a).)
\item
  There exist positive measures \(\alpha\) and \(\beta\) on X such that \(f(x)-cf(x')=\alpha(f)\) and \(f(x')-cf(x)=\beta(f)\) for \(f\in R\) . Setting \(\mu_{x}=(1-c^{2})^{-1}(c\beta+\alpha)\) , \(\mu_{x'}=(1-c^{2})^{-1}(c\alpha+\beta)\) , one has \(\mu_{x}(f)=f(x)\) and \(\mu_{x'}(f)=f(x')\) for all \(f\in B\) , and \(c\mu_{x}\leqslant\mu_{x'}\) , \(c\mu_{x'}\leqslant\mu_{x}\) .
\item
  We still assume \(x \sim x'\). Show that if \(\mu\) is a positive measure on X such that \(\mu(f) = f(x)\) for all \(f \in B\), there exists a positive measure \(\mu'\) on X such that \(\mu'(f) = f(x')\) for all \(f \in B\) and \(c\mu \leqslant \mu'\) for some \(c \in]0,1[\). (With the notations of c), set \(\mu' = \mu_{x'} - c\mu_x + c\mu\).)
\item
  The relation \(x \sim x'\) is an equivalence relation on X. (Use d).)
\end{enumerate}

¶ 10) Let X be a compact topological space and A a unital Banach subalgebra of \(\mathcal{C}(\mathrm{X})\) for the induced norm.

\begin{enumerate}
\def\labelenumi{\alph{enumi})}
\item
  A subset K of X is said to be antisymmetric (relative to A) if every function \(f \in A\) which is real on K is constant on K. Every antisymmetric subset is contained in a maximal antisymmetric subset. A maximal antisymmetric subset is closed. The set \(\mathfrak{R}\) of maximal antisymmetric subsets is a partition of X.
\item
  Let \(A^{\perp}\) be the orthogonal of A in the Banach space of complex measures on X. Let \(\mu\) be an extreme point of the unit ball of \(A^{\perp}\). Show that the support of \(\mu\) is antisymmetric.
\item
  Let \(f \in \mathcal{C}(\mathrm{X})\). If \(f|K \in A|K\) for all \(K \in \mathfrak{R}\), then \(f \in A\). (Apply b) and the Krein–Milman theorem.) From this, recover the Weierstrass–Stone theorem (TG, X, p.~36, th. 3).
\end{enumerate}

\(d)\) A nonempty subset E of X is called a peak (relative to A) if there exists \(f\in \mathrm{A}\) such that \(\| f\| = 1\) and E is the set of those \(x\in \mathrm{X}\) for which \(f(x) = 1\). One may then assume \(|f(x)| < 1\) for \(x\notin \mathrm{E}\) (replace \(f\) by \(\frac{1}{2} (1 + f)\)). Every nonempty countable intersection of peaks is a peak. If \(\mathrm{E}_1,\mathrm{E}_2\) are peaks, then \(\mathrm{E}_1\cup \mathrm{E}_2\) is a peak. (Let \(f_{1},f_{2}\in \mathrm{A}\) with \(f_{i} = 1\) on \(\mathrm{E}_i\), \(|f_i| < 1\) on \(\mathrm{X - E}_i\); let \(g_{i} = \frac{1}{4} (1 - f_{i})^{1 / 3}\in \mathrm{A}\); then \(1 - g_{1}g_{2} = 1\) on \(\mathrm{E}_1\cup \mathrm{E}_2\), \(|1 - g_1g_2| < 1\) on \(\mathrm{X - (E_1\cup E_2)}\).)

\begin{enumerate}
\def\labelenumi{\alph{enumi})}
\setcounter{enumi}{4}
\item
  Let E be an intersection of peaks. Let I be the set of \(f \in A\) vanishing on E. Then \(f \mapsto f|E\) defines an isometry from A/I onto the algebra A\textbar E of restrictions to E of functions in A.
\item
  Let E be a peak. Let \(E' \subset E\) be a peak relative to A\textbar E. Then \(E'\) is a peak relative to A.
\item
  Every antisymmetric subset of X is an intersection of peaks. (Use d) and f).)
\item
  For every antisymmetric subset K of X, the algebra A\textbar K is a closed subalgebra of \(\mathcal{C}(\mathrm{K})\). (Use \(e\) and \(g\).)
\item
  Let R be the set of real parts of elements of A. Suppose that \(X \in \mathfrak{R}\) and that R is stable under multiplication. Then X consists of a single point. (Let \(x_{0} \in X\). For \(u \in R\), there exists a unique \(f \in A\) such that \(u = \mathcal{R}(f)\) and \(f(x_{0}) \in \mathbf{R}\); set \(\mathrm{N}(u) = \|f\|\). Then R is a real Banach space for N; the closed graph theorem proves that multiplication in R is separately continuous, hence continuous. Deduce from this a Banach algebra structure on \(S = R + iR\). Let \(p \in R\) with \(p(x) > 0\) for all \(x \in X\). Show that, for every character \(\chi\) of S, one has \(\mathcal{R}(\chi(p)) > 0\), and deduce that \(\log(p) \in R\). Assuming \(X \neq \{x_{0}\}\), one can construct an \(f \in A\) such that \(0 < \mathcal{R}(f) \leqslant 1\) on X, \(f(x_{0}) \in \mathbf{R}\), and such that \(\|f\|\) is arbitrarily large. According to what precedes, there exists \(V \in A\) with \(|V|^{2} = \mathcal{R}(f)\). We have \(\|V\| \leqslant 1\), hence \(\mathrm{N}(\mathcal{R}(V)) \leqslant 2\), but
\end{enumerate}

\[
\mathrm{N} (\mathcal {R} (\mathrm{V}) ^ {2}) \geqslant \frac {1}{2} \left(\| \mathrm{V} ^ {2} + f \| - | \mathcal {I} (\mathrm{V}) (x _ {0}) ^ {2} |\right)
\]

is arbitrarily large.)

\begin{enumerate}
\def\labelenumi{\alph{enumi})}
\setcounter{enumi}{9}
\tightlist
\item
  If R is stable under multiplication, then A = C(X). (Use a), c), h), i).) Consequently, if A \(\neq\) C(X), there exists \(u \in R\) such that \(u^{2} \notin R\).
\end{enumerate}

\begin{enumerate}
\def\labelenumi{\arabic{enumi})}
\setcounter{enumi}{10}
\tightlist
\item
  Let X be a metric space, whose distance is denoted by d. A function \(f: X \to C\) is said to be Lipschitz if there exists a real number \(c \geqslant 0\) such that, for all \(x\) and \(y\) in X, one has
\end{enumerate}

\[
| f (x) - f (y) | \leqslant c d (x, y)
\]

(cf.~FVR, IV, p.~7, def. 1). We denote by \(\mathrm{Lip}_b(\mathrm{X})\) the set of bounded Lipschitz functions on X.

\begin{enumerate}
\def\labelenumi{\alph{enumi})}
\item
  The set \(\mathrm{Lip}_b(\mathrm{X})\) is a subalgebra of \(\mathcal{C}(\mathrm{X})\).
\item
  Equipped with the norm
\end{enumerate}

\[
\| f\| = \sup_{x\in \mathrm{X}}|f(x)| + \sup_{\substack{x,y\in \mathrm{X}\\ x\neq y}}\frac{|f(x) - f(y)|}{d(x,y)},
\]

and with the involution \(f \mapsto \overline{f}\), the algebra \(\mathrm{Lip}_b(\mathrm{X})\) is a commutative involutive Banach algebra. It is not a stellar algebra in general.

\begin{enumerate}
\def\labelenumi{\alph{enumi})}
\setcounter{enumi}{2}
\item
  For any commutative Banach algebra A, we denote by \(\mathsf{X}_{m}(\mathsf{A})\) the metric space \(\mathsf{X}(\mathsf{A})\) endowed with the distance induced by the norm of the dual \(A'\) of A. The metric space \(\mathsf{X}_{m}(\mathsf{A})\) is complete and of diameter \(\leqslant 2\) .
\item
  Let A be a unital commutative Banach algebra. The Gelfand transform maps A into \(\operatorname{Lip}_{b}(\mathsf{X}_{m}(\mathsf{A}))\). If A is semisimple, then it induces a homeomorphism from A onto its image.
\end{enumerate}

In the sequel, we assume that X is a compact metric space. We denote by \(A = \operatorname{Lip}_{b}(X)\).

\begin{enumerate}
\def\labelenumi{\alph{enumi})}
\setcounter{enumi}{4}
\item
  The map \(\varphi\) which sends \(x \in X\) to the character \(f \mapsto f(x)\) is a homeomorphism from \(X\) onto \(X(A)\).
\item
  Let \(d_{1}\) be the distance on X induced by the homeomorphism \(\varphi\). Then \(d_{1}\) is equivalent to d.
\end{enumerate}

\chapter{Locally compact abelian groups}

Throughout this chapter, the letter G denotes, unless otherwise stated, a locally compact abelian group equipped with a Haar measure generally denoted by dx; for \(p \in [1, +\infty]\), the space \(L^{p}(G, dx)\) will simply be denoted \(L^{p}(G)\), and its norm will be denoted by \(f \mapsto \|f\|_{p}\). We identify \(L^{1}(G)\) with a subspace of \(\mathcal{M}^{1}(G)\) by the map \(f \mapsto f \cdot dx\). Recall that the support of the Haar measure is equal to G (INT, VII, §1, no. 1, remark 3); in particular (INT, III, §2, no. 2, proposition 9), the canonical mapping from the space \(\mathcal{K}(G)\) into \(L^{p}(G)\) is injective for \(p \in [1, +\infty]\). For \(p \neq +\infty\), we identify \(L^{p}(G)\) with a subspace of the space \(\widetilde{\mathcal{F}}(G; \mathbf{C})\), the space of classes of complex-valued functions on G defined and finite almost everywhere (INT, IV, §3, nos. 5, 6). In particular, the notation \(L^{1}(G) \cap L^{2}(G)\) denotes the intersection of \(L^{1}(G)\) and \(L^{2}(G)\) in this space. We denote by \(f \mapsto \widetilde{f}\) the involution on the involutive algebra \(L^{1}(G)\) (example 4 of I, p.~99); we have \(\widetilde{f}(x) = \overline{f(x^{-1})}\) for all \(x\) in G.

Recall (INT, V, §5, no. 3, th. 1) that if μ is a complex measure on a locally compact topological space X and if f is a locally μ-integrable function on X, then the measure ν with density f with respect to μ is denoted f·μ or simply fμ. A function g from X to C is essentially integrable for the measure ν if and only if gf is essentially integrable for \(\mu\); we then have

\[
(f \cdot \mu) (g) = \int_ {\mathrm{X}} g d (f \cdot \mu) = \int_ {\mathrm{X}} g f d \mu = \mu (g f).
\]

If G is a compact topological group, the normalized Haar measure on G is the unique Haar measure \(\mu\) on G such that \(\mu(\mathrm{G}) = 1\).

If \(X\) is a discrete topological space, the counting measure on \(X\) is the discrete measure \(\mu\) on \(X\) such that \(\mu(\{x\}) = 1\) for all \(x \in X\). If \(G\) is a discrete topological group, the counting measure on \(G\) is a Haar measure on \(G\).

We will also make use of the following two lemmas.

Lemma 1. --- Let X be a locally compact topological space and \(\mu\) a positive measure on X. Let \(x\) be an element of the support of \(\mu\). Let U be an open neighborhood of \(x\). There exists a function \(f \in \mathcal{K}_{+}(\mathrm{X})\) whose support is contained in U such that \(\int fd\mu = 1\).

By definition of the support of a measure (INT, III, §2, no. 2, def. 1), there exists a function \(g \in \mathcal{K}(\mathrm{X})\) supported in U such that \(\mu(g) \neq 0\) . We then have \(\mu(|g|) > 0\) since \(\mu\) is positive, and the function \(f = \mu(|g|)^{-1}|g|\) has the desired properties.

Lemma 2. --- Let G and H be topological groups with identity elements \(e_{\mathrm{G}}\) and \(e_{\mathrm{H}}\). Suppose the topological group G is Hausdorff. Let f be a morphism of topological groups from G to H. Suppose that for every neighborhood U of \(e_{\mathrm{G}}\) in G, there exists a neighborhood W of \(e_{\mathrm{H}}\) in H such that \(f^{-1}(\mathrm{W}) \subset \mathrm{U}\). Then the morphism f is injective and strict (TG, III, p.~16, def. 1).

The hypotheses imply that \(\operatorname{Ker}(f)\) is contained in every neighborhood of \(e_{G}\) in G, hence that \(\operatorname{Ker}(f) = \{e_{\mathrm{G}}\}\) since G is separated. The homomorphism from G to \(f(\mathrm{G})\) induced by f by passing to subspaces is then bijective; denote by \(g: f(\mathrm{G}) \to \mathrm{G}\) the inverse homomorphism. The hypotheses then imply that g is continuous at \(e_{H}\), therefore continuous (TG, III, p.~15, prop. 23).

\section{FOURIER TRANSFORMATION}

\subsection{Unitary characters of a locally compact abelian group}

DEFINITION 1. --- A unitary character of G is a continuous homomorphism from G into the multiplicative group U of complex numbers of modulus 1.

In other words, a unitary character is a continuous function \(\chi\) on G, with complex values, such that:

\[
\chi (x y) = \chi (x) \chi (y), \qquad | \chi (x) | = 1 \qquad (x, y \in \mathrm{G}).
\]

In this chapter, we will often simply say “character” instead of “unitary character.”

Let E be a Hilbert space of dimension 1, and let \(\chi\) be a unitary character of G. The map that assigns to \(x \in \mathrm{G}\) the homothety of ratio \(\chi(x)\) in E is a continuous linear isometric representation of G on E. Conversely, every bounded continuous linear representation of G on E is obtained by this process, and in particular is unitary.

It is immediate that the product of two unitary characters, the inverse of a unitary character, and the constant function equal to 1 are unitary characters. Consequently, the set \(\widehat{\mathbf{G}}\) of unitary characters of G forms a group under multiplication. This group is commutative. On the other hand, the map \((\chi_1,\chi_2)\mapsto \chi_1\chi_2^{-1} = \chi_1\overline{\chi}_2\) is continuous for the compact-open topology, and \(\widehat{\mathbf{G}}\), endowed with the topology of compact convergence, is a topological group (TG, X, p.~6, corollary 2 and remark 1).

DEFINITION 2. --- The topological group \(\widehat{\mathbf{G}}\) is called the dual group of G.

Since G is locally compact, the map \((x,\chi)\mapsto\chi(x)\) is continuous on \(G\times\widehat{G}\) (TG, X, p.~28, th. 3).

Recall that \(\mathcal{M}^{1}(G)\) denotes the unital involutive Banach algebra of bounded complex measures on G (Example 4 of I, p.~99). For every bounded complex measure \(\mu\in\mathcal{M}^{1}(G)\) and every \(\chi\in\widehat{G}\) , we denote by

\[
\chi (\mu) = \int_ {\mathrm{G}} \chi (x) d \mu (x)\tag{1}
\]

(cf. INT, VIII, §2, no. 6).

Lemma 1. --- For every \(\chi \in \widehat{\mathbf{G}}\), the map \(\mu \mapsto \chi(\mu)\) is a Hermitian character of the involutive Banach algebra \(\mathcal{M}^1(\mathbf{G})\).

By INT, VIII, §3, no. 3, prop. 11, the map \(\mu \mapsto \chi(\mu)\) is a character of the involutive Banach algebra \(\mathcal{M}^{1}(G)\). Moreover, we have:

\[
\chi (\mu^ {*}) = \int_ {\mathrm{G}} \chi (x ^ {- 1}) d \overline {{\mu}} (x) = \int_ {\mathrm{G}} \overline {{\chi (x)}} d \overline {{\mu}} (x) = \overline {{\chi (\mu)}}
\]

and this character is therefore Hermitian.

Thus we have defined a mapping from \(\widehat{G}\) into \(\mathsf{X}(\mathcal{M}^{1}(\mathrm{G}))\); it will be called canonical.

Let \(\chi \in \widehat{\mathbf{G}}\). The restriction of \(\mu \mapsto \chi(\mu)\) to \(\mathrm{L}^1(\mathrm{G})\) is nonzero (cf. INT, VIII, §2, no. 7, prop. 10). By restriction to the involutive Banach subalgebra \(\mathrm{L}^1(\mathrm{G})\), we therefore obtain a map from \(\widehat{\mathbf{G}}\) into \(\mathsf{X}(\mathrm{L}^1(\mathrm{G}))\), said to be canonical. It associates to \(\chi \in \widehat{\mathbf{G}}\) the Hermitian character

\[
f \mapsto \chi (f) = \chi (f \cdot d x) = \int_ {\mathrm{G}} f (x) \chi (x) d x \qquad (f \in \mathrm{L} ^ {1} (\mathrm{G}))\tag{2}
\]

of \(L^{1}(G)\) .

PROPOSITION 1. --- The canonical map from \(\widehat{\mathbf{G}}\) into \(\mathsf{X}(\mathrm{L}^{1}(\mathbf{G}))\) is a homeomorphism.

Let us denote this canonical map by ev and, for \(\chi \in \widehat{\mathbf{G}}\), let \(\mathrm{ev}_{\chi}\) be the image of the character \(\chi\) by ev, that is, the Hermitian character \(f \mapsto \chi(f)\) of \(\mathrm{L}^1(\mathrm{G})\). Considered as a mapping from \(\widehat{\mathbf{G}}\) into the dual of \(\mathrm{L}^1(\mathrm{G})\), the map ev is the composition of the injection of \(\widehat{\mathbf{G}}\) into \(\mathrm{L}^{\infty}(\mathrm{G})\), equipped with the weak topology \(\sigma(\mathrm{L}^{\infty}(\mathrm{G}), \mathrm{L}^{1}(\mathrm{G}))\), and the injection of \(\mathrm{L}^{\infty}(\mathrm{G})\) into the dual of \(\mathrm{L}^1(\mathrm{G})\), equipped with the topology of pointwise convergence. Since \(\widehat{\mathbf{G}}\) is a bounded subset of \(\mathrm{L}^{\infty}(\mathrm{G})\), the first map is continuous by Lebesgue's theorem (INT, IV, §3, no. 7, th. 6). The second is also continuous, by definition. This proves that the map ev is continuous.

If \(\chi \in \widehat{\mathrm{G}}\) and if \(f \in \mathrm{L}^1(\mathrm{G})\) , one has \(\mathrm{ev}_{\chi}(\varepsilon_x * f) = \chi(x)\mathrm{ev}_{\chi}(f)\) . Taking \(f\) such that \(\mathrm{ev}_{\chi}(f) \neq 0\) , we deduce that the map ev is injective.

Let \(\zeta\in\mathsf{X}(\mathrm{L}^{1}(\mathrm{G}))\) and let \(f\in\mathrm{L}^{1}(\mathrm{G})\) such that \(\zeta(f)\neq0\) . Let us define a map \(\chi\colon\mathrm{G}\to\mathbf{C}\) by setting, for \(x\in\mathrm{G}\) :

\[
\chi (x) = \frac {\zeta (\varepsilon_ {x} * f)}{\zeta (f)}.\tag{3}
\]

We have \(\chi(e) = 1\). Since the map \(x \mapsto \varepsilon_x * f = \gamma(x)(f)\) from G to L \(^1\) (G) is continuous (INT, VIII, §2, n° 5, prop. 8), the map \(\chi\) is continuous. It is bounded because, for all \(x \in \mathbf{G}\), one has

\[
| \chi (x) | \leqslant \frac {\| \varepsilon_ {x} * f \|}{| \zeta (f) |} = \frac {\| f \|}{| \zeta (f) |}
\]

(th. 1 of I, p.~29 and INT, VIII, loc. cit.).

Now let \(\mathfrak{B}\) be a base of the filter of neighborhoods of \(e\) consisting of compact neighborhoods. For every \(V \in \mathfrak{B}\), let \(g_V\) be a positive continuous function, vanishing outside \(V\) and with integral equal to 1 (lemma 1 of II, p.~200). For every function \(h \in L^1(G)\), we then have

\[
\varepsilon_ {x} * h = \lim _ {\mathrm{V}, \mathfrak {B}} (\varepsilon_ {x} * h) * g _ {\mathrm{V}} = \lim _ {\mathrm{V}, \mathfrak {B}} (\varepsilon_ {x} * g _ {\mathrm{V}} * h),
\]

in L \(^{1}\) (G), the limit being taken along the filter of sections of \(\mathfrak{B}\) (INT, VIII, §4, n° 7, prop. 20). In particular, since \(\zeta(\varepsilon_x * g_V * f) = \zeta(\varepsilon_x * g_V)\zeta(f)\), one deduces that

\[
\chi (x) = \lim _ {\mathrm{V}, \mathfrak {B}} \zeta (\varepsilon_ {x} * g _ {\mathrm{V}}).
\]

For every \(h\in \mathrm{L}^1 (\mathrm{G})\) , we obtain

\[
\zeta (\varepsilon_ {x} * h) = \lim _ {\mathrm{V}, \mathfrak {B}} \zeta (\varepsilon_ {x} * g _ {\mathrm{V}} * h) = \zeta (h) \lim _ {\mathrm{V}, \mathfrak {B}} \zeta (\varepsilon_ {x} * g _ {\mathrm{V}}) = \chi (x) \zeta (h).
\]

Consequently, we have, for \(x, y \in G\):

\[
\chi (x y) = \frac {\zeta (\varepsilon_ {x} * \varepsilon_ {y} * f)}{\zeta (f)} = \frac {\chi (x) \zeta (\varepsilon_ {y} * f)}{\zeta (f)} = \chi (x) \chi (y),
\]

which proves that \(\chi\) is a homomorphism from G to \(\mathbf{C}^*\). Since \(\chi\) is bounded and continuous, it is a unitary character of G. Moreover, if \(g \in \mathrm{L}^{1}(\mathrm{G})\), one has

\[
g * f = \int_ {\mathrm{G}} (\varepsilon_ {x} * f) g (x) d x
\]

in L \(^{1}\) (G) (INT, VIII, §1, n° 5, prop. 7), whence

\[
\begin{array}{c} \zeta (g) \zeta (f) = \zeta (g * f) = \int_ {\mathrm{G}} \zeta (\varepsilon_ {x} * f) g (x) d x \\ = \zeta (f) \int_ {\mathrm{G}} \chi (x) g (x) d x = \mathrm{ev} _ {\chi} (g) \zeta (f) \end{array}
\]

(INT, VI, §1, n°1, prop. 1), which shows that \(\zeta = \mathrm{ev}_{\chi}\) . Consequently, ev is surjective, hence bijective.

Finally let us show that the inverse map \(ev^{-1}\) is continuous. Let \(\zeta \in \mathsf{X}(\mathrm{L}^{1}(\mathrm{G}))\). Let \(f \in \mathrm{L}^{1}(\mathrm{G})\) be a function such that \(\zeta(f) \neq 0\). The set W of \(\xi \in \mathsf{X}(\mathrm{L}^{1}(\mathrm{G}))\) such that \(\xi(f) \neq 0\) is an open neighborhood of \(\zeta\) in \(\mathsf{X}(\mathrm{L}^{1}(\mathrm{G}))\). For every \(\xi \in W\), what precedes shows that \(\mathrm{ev}^{-1}(\xi)\) is the character

\[
x \mapsto \frac {\xi (\varepsilon_ {x} * f)}{\xi (f)}.
\]

Let \(\mathfrak{F}\) be a filter on \(W \subset X(L^{1}(G))\) converging to \(\zeta\). Since the set \(X(L^{1}(G))\) is bounded, hence equicontinuous, in \(L^{\infty}(G)\), the uniform structure of simple convergence coincides with the uniform structure of compact convergence (TG, X, p.~16, th. 1). Let K be a compact subset of G. The set of \(\varepsilon_{x} * f\) for \(x \in K\) is compact in \(L^{1}(G)\) (INT, VIII, §2, n° 5, prop. 8). We therefore have

\[
\lim _ {\xi , \mathfrak {F}} \xi (\varepsilon_ {x} * f) = \zeta (\varepsilon_ {x} * f)
\]

uniformly for \(x \in K\) . From this we deduce that

\[
\lim _ {\xi , \mathfrak {F}} \mathrm{ev} ^ {- 1} (\xi) = \mathrm{ev} ^ {- 1} (\zeta),
\]

hence ev \(^{-1}\) is continuous at \(\zeta\) . This completes the proof of the proposition.

From now on we identify a unitary character \(\chi\) of G with the character \(f\mapsto\int_{\mathrm{G}}f(x)\chi(x)dx\) of \(\mathrm{L}^{1}(\mathrm{G})\).

Remarks. --- 1) *The bijectivity of the map from \(\widehat{G}\) onto \(X(L^{1}(G))\) in prop. 1 is a special case of the correspondence between continuous representations of a locally compact group H (not necessarily commutative) and continuous representations of the algebra \(L^{1}(H)\).*

\begin{enumerate}
\def\labelenumi{\arabic{enumi})}
\setcounter{enumi}{1}
\tightlist
\item
  The canonical map from \(\widehat{G}\) into \(\mathsf{X}(\mathcal{M}^{1}(\mathsf{G}))\) is not surjective in general (II, p.~308, exerc. 14).
\end{enumerate}

COROLLARY 1. --- Every character of L \(^{1}\) (G) is Hermitian. The canonical map from X(Stell(G)) (def. 9 of I, p.~125) into X(L \(^{1}\) (G)) is a homeomorphism.

The first assertion follows from Proposition 1 and Lemma 1, restricted to \(L^{1}(G)\). The second follows from the first and from the corollary of Prop. 20 of I, p.~124.

Corollary 2. --- The topological group \(\widehat{\mathbf{G}}\) is locally compact.

Indeed, \(\mathsf{X}(\mathrm{L}^{1}(\mathrm{G}))\) is locally compact (corollary of Theorem 1 of I, p.~29).

We shall identify \(\widehat{G}\) with \(\mathsf{X}(\mathrm{L}^1 (\mathrm{G}))\) and \(\mathsf{X}(\mathrm{Stell}(\mathrm{G}))\). For \(x\in G\) and \(\chi \in \widehat{G}\), we shall denote by \(\langle \chi ,x\rangle\) the complex number \(\chi (x)\), which belongs to \(\mathbf{U}\).

We say that \(x\) and \(\chi\) are orthogonal if \(\langle \chi, x \rangle = 1\). Let A be a subset of G (resp. of \(\widehat{\mathbf{G}}\)); the set of elements of \(\widehat{\mathbf{G}}\) (resp. of G) orthogonal to A is a closed subgroup of \(\widehat{\mathbf{G}}\) (resp. of G) which is called the orthogonal of A and is denoted \(\mathrm{A}^{\perp}\). The orthogonal of G consists only of \(e\).

For \(x \in G\), let \(\eta(x)\) be the map from \(\widehat{G}\) to U defined by \(\chi \mapsto \langle \chi, x \rangle\). By definition of multiplication in \(\widehat{G}\), the map \(\eta(x)\) is a group homomorphism. It is continuous since the map \((x, \chi) \mapsto \langle \chi, x \rangle\) from \(G \times \widehat{G}\) to U is continuous (TG, X, p.~28, th. 3). We have thus defined a mapping \(\eta\), called canonical, from G into the bidual group \(\widehat{\widehat{G}}\); it is a group homomorphism. Moreover, the map \(\eta\) is continuous (TG, X, p.~28, th. 3). We shall prove later (II, p.~220, th. 2) that \(\eta\) is a topological group isomorphism from G onto \(\widehat{\widehat{G}}\).

Let G and H be locally compact abelian groups, and \(\varphi\colon\mathrm{G}\to\mathrm{H}\) a morphism of topological groups. For every \(\chi\in\widehat{\mathrm{H}}\), the map \(\chi\circ\varphi\) is a character of G, denoted \(\widehat{\varphi}(\chi)\). This definition translates to the formula

\[
\langle \chi , \varphi (x) \rangle = \langle \widehat {\varphi} (\chi), x \rangle\tag{4}
\]

for all \(\chi\in\widehat{H}\) and \(x\in G\). From this we deduce that \(\widehat{\varphi}\) is a morphism from the topological group \(\widehat{H}\) to the topological group \(\widehat{G}\); we say that \(\widehat{\varphi}\) is the dual of the morphism \(\varphi\).

Let K be a locally compact commutative group and \(\psi\colon\mathrm{H}\to\mathrm{K}\) a morphism of topological groups. The definition shows that \(\widehat{\psi\circ\varphi}=\widehat{\varphi}\circ\widehat{\psi}\). If \(\varphi\) is the identity map of G, then \(\widehat{\varphi}\) is the identity map of \(\widehat{\mathrm{G}}\). In particular, if \(\varphi\) is an isomorphism of topological groups, then so is \(\widehat{\varphi}\), and \(\widehat{\varphi}^{-1}\) is the dual of \(\varphi^{-1}\).

Lemma 2. --- Let G and H be locally compact abelian groups, and let \(f\colon\mathrm{H}\to\mathrm{G}\) be a morphism of topological groups. The kernel of \(\widehat{f}\) is the orthogonal of the image of \(f\).

By definition, we have \(\chi \in \mathrm{Ker}(\widehat{f})\) if and only if the restriction of \(\chi\) to the image of \(f\) is trivial.

PROPOSITION 2. --- Let \(n \geqslant 0\) be an integer and let \(\mathrm{G}_{1}, \ldots, \mathrm{G}_{n}\) be locally compact commutative groups. Let G be the product group of the groups \(\mathrm{G}_{j}\) for \(1 \leqslant j \leqslant n\). For \(1 \leqslant j \leqslant n\), let \(\lambda_{j}\) be the injection of \(\mathrm{G}_{j}\) into G which associates to \(x \in \mathrm{G}_{j}\) the element \((x_{k})\) such that \(x_{k} = e\) if \(k \neq j\) and \(x_{j} = x\). The map

\[
(\widehat {\lambda} _ {j}) _ {1 \leqslant j \leqslant n} \colon \widehat {\mathrm{G}} \to \prod_ {1 \leqslant j \leqslant n} \widehat {\mathrm{G}} _ {j}
\]

is an isomorphism of topological groups.

Let \(m\) be the product map \(\widehat{\mathbf{G}}^n\) to \(\widehat{\mathbf{G}}\), and, for every \(j\) such that \(1 \leqslant j \leqslant n\), let \(\pi_j\) be the projection of \(\mathbf{G}\) onto \(\mathbf{G}_j\). Let \(\mu\) be the morphism of topological groups \(m \circ (\widehat{\pi}_j)_j\) from \(\prod \widehat{\mathbf{G}}_j\) to \(\widehat{\mathbf{G}}\), so that

\[
\langle \mu ((\chi_ {j})), (x _ {j}) \rangle = \prod_ {j = 1} ^ {n} \langle \chi_ {j}, x _ {j} \rangle .
\]

The map \(\mu\) is continuous, and one verifies that \(\mu\) and \((\widehat{\lambda}_j)\) are inverse bijections of each other. The proposition follows.

Remark. --- The computation of the dual group of an infinite product of compact commutative groups is the subject of Corollary 4 of II, p.~234 below. The case of a locally compact commutative group that is an arbitrary product of locally compact groups follows from these two statements, since in such a product all but finitely many factors are compact (TG, I, p.~66, Prop. 14, b)).

\subsection{Definition of the Fourier transform}

DEFINITION 3. --- Let \(\mu \in \mathcal{M}^{1}(\mathrm{G})\) be a bounded complex measure on G. The Fourier transform of \(\mu\) is the function \(\mathcal{F}_{\mathrm{G}}(\mu)\) on \(\widehat{\mathrm{G}}\) defined by

\[
\mathcal {F} _ {\mathrm{G}} (\mu) (\widehat {x}) = \int_ {\mathrm{G}} \overline {{\langle \widehat {x} , x \rangle}} d \mu (x).\tag{5}
\]

We call the Fourier cotransform of \(\mu\) the function \(\overline{\mathcal{F}}_{\mathrm{G}}(\mu)\) on \(\widehat{\mathrm{G}}\) defined by

\[
\overline {{\mathcal {F}}} _ {\mathrm{G}} (\mu) (\widehat {x}) = \int_ {\mathrm{G}} \langle \widehat {x}, x \rangle d \mu (x).\tag{6}
\]

When there is no ambiguity regarding the group G under consideration, we will also write \(\mathcal{F}(\mu)\) and \(\overline{\mathcal{F}}(\mu)\). We also sometimes denote by \(\widehat{\mu} = \mathcal{F}_{\mathrm{G}}(\mu)\).

PROPOSITION 3. --- For every measure \(\mu \in \mathcal{M}^{1}(\mathrm{G})\), the functions \(\mathcal{F}_{\mathrm{G}}(\mu)\) and \(\overline{\mathcal{F}}_{\mathrm{G}}(\mu)\) are continuous and bounded. The maps \(\mathcal{F}_{\mathrm{G}}: \mu \mapsto \mathcal{F}_{\mathrm{G}}(\mu)\) and \(\overline{\mathcal{F}}_{\mathrm{G}}: \mu \mapsto \overline{\mathcal{F}}_{\mathrm{G}}(\mu)\) are continuous morphisms of the involutive algebra \(\mathcal{M}^{1}(\mathrm{G})\) into the involutive algebra of bounded continuous functions on \(\widehat{\mathrm{G}}\) (example 1 of I, p.~99).

Let \(\mu \in \mathcal{M}^1 (\mathrm{G})\). For every \(\chi \in \widehat{\mathrm{G}}\), one has

\[
| \mathcal {F} (\mu) (\chi) | = \left| \int_ {\mathrm{G}} \overline {{\langle \chi , x \rangle}} d \mu (x) \right| \leqslant \| \mu \| _ {1},\tag{7}
\]

so the Fourier transform of \(\mu\) is bounded. Similarly, one verifies that \(\overline{\mathcal{F}}(\mu)\) is bounded.

If \(\chi\) tends to \(\chi_0\) in \(\widehat{\mathbf{G}}\), the function \(\chi\) on \(\mathbf{G}\) tends to \(\chi_0\) uniformly on every compact set while remaining bounded by the constant function 1, which belongs to \(\mathrm{L}^1 (\mathrm{G},\mu)\). By Lebesgue's theorem (INT, IV, §3, n°7, th. 6), it follows that \(\mathcal{F}(\mu)(\chi)\) tends to \(\mathcal{F}(\mu)(\chi_0)\). Therefore \(\mathcal{F}(\mu)\) is continuous. The same is true of \(\overline{\mathcal{F}} (\mu)\).

For every \(\chi \in \widehat{\mathbf{G}}\), the map \(\mu \mapsto \chi(\mu) = \int \langle \chi, x \rangle d\mu(x)\) is a Hermitian character of \(\mathcal{M}^1(\mathbf{G})\) (lemma 1 of II, p.~202). This implies that \(\mathcal{F}\) and \(\overline{\mathcal{F}}\) are involutive algebra morphisms of \(\mathcal{M}^1(\mathbf{G})\) into the involutive algebra of bounded continuous functions on \(\widehat{\mathbf{G}}\). The inequality (7) shows that these morphisms are continuous.

The Fourier transform of G (resp. the Fourier cotransform of G) is the mapping \(\mu \mapsto \mathcal{F}_{\mathrm{G}}(\mu)\) (resp. the map \(\mu \mapsto \overline{\mathcal{F}}_{\mathrm{G}}(\mu)\)) from \(\mathcal{M}^1 (\mathrm{G})\) into \(\mathcal{C}_b(\widehat{\mathrm{G}})\).

Let us note some useful formulas for \(\mu\in\mathcal{M}^{1}(G)\), \(x\in G\), and \(\chi\in\widehat{G}\):

\begin{enumerate}
\def\labelenumi{(\arabic{enumi})}
\setcounter{enumi}{7}
\tightlist
\item
\end{enumerate}

\[
\overline {{\mathcal {F}}} (\mu) (\chi) = \mathcal {F} (\mu) (\chi^ {- 1}) = \overline {{\mathcal {F} (\overline {{\mu}}) (\chi)}},\tag{9}
\]

\[
\| \mathcal {F} (\mu) \| _ {\infty} = \| \overline {{\mathcal {F}}} (\mu) \| _ {\infty} \leqslant \| \mu \| _ {1},\tag{10}
\]

\[
\left\{ \begin{array}{l} \mathcal {F} (\varepsilon_ {x}) (\chi) = \overline {{\langle \chi , x \rangle}}, \\ \overline {{\mathcal {F}}} (\varepsilon_ {x}) (\chi) = \langle \chi , x \rangle , \end{array} \right.
\]

(in particular \(\mathcal{F}(\varepsilon_e) = \overline{\mathcal{F}} (\varepsilon_e) = 1\) ),

\begin{enumerate}
\def\labelenumi{(\arabic{enumi})}
\setcounter{enumi}{10}
\tightlist
\item
\end{enumerate}

\[
\left\{ \begin{array}{l} \mathcal {F} (\varepsilon_ {x} * \mu) (\chi) = \overline {{\langle \chi , x \rangle}} \mathcal {F} (\mu) (\chi), \\ \overline {{\mathcal {F}}} (\varepsilon_ {x} * \mu) (\chi) = \langle \chi , x \rangle \overline {{\mathcal {F}}} (\mu) (\chi), \end{array} \right.\tag{12}
\]

\[
\left\{ \begin{array}{l} \mathcal {F} (\chi \cdot \mu) = \varepsilon_ {\chi} * \mathcal {F} (\mu), \\ \overline {{\mathcal {F}}} (\chi \cdot \mu) = \varepsilon_ {\chi^ {- 1}} * \overline {{\mathcal {F}}} (\mu). \end{array} \right.
\]

Formulas (8), (9), (10), and (11) follow from the definitions. Let us prove the first of formulas (12), the second being analogous. For all \(\xi\) in \(\widehat{G}\), the equalities

\[
\begin{array}{c} \mathcal {F} (\chi \cdot \mu) (\xi) = \int_ {\mathrm{G}} \overline {{\langle \xi , x \rangle}} \langle \chi , x \rangle d \mu (x) = \int_ {\mathrm{G}} \overline {{\langle \xi \chi^ {- 1} , x \rangle}} d \mu (x) \\ = \mathcal {F} (\mu) (\xi \chi^ {- 1}) = (\varepsilon_ {\chi} * \mathcal {F} (\mu)) (\xi). \end{array}
\]

Note moreover that for all \(\chi\in\widehat{\mathrm{G}}\) and all measures \(\mu\) and \(\nu\) in \(\mathcal{M}^{1}(\mathrm{G})\), one has

\[
(\varepsilon_ {\chi} * \mathcal {F} (\mu)) (\varepsilon_ {\chi} * \mathcal {F} (\nu)) = \varepsilon_ {\chi} * (\mathcal {F} (\mu) \mathcal {F} (\nu)),\tag{13}
\]

since both sides of this equality are functions on \(\widehat{G}\) whose value at \(\xi\in\widehat{G}\) is

\[
\mathcal {F} (\mu) (\xi \chi^ {- 1}) \mathcal {F} (\nu) (\xi \chi^ {- 1}).
\]

Let H be a locally compact abelian group, and let \(\varphi\colon\mathrm{G}\to\mathrm{H}\) be a continuous morphism. Let \(\mu\in\mathcal{M}^{1}(\mathrm{G})\). The image measure \(\varphi(\mu)\) is defined (INT, V, §6, no. 1, remark 1), and it follows that \(\mathcal{F}_{\mathrm{H}}(\varphi(\mu))=\mathcal{F}_{\mathrm{G}}(\mu)\circ\widehat{\varphi}\) (cf. INT, V, §6, no. 4, prop. 7).

By restriction to the subalgebra \(L^{1}(G)\) of \(\mathcal{M}^{1}(G)\), we obtain the definitions of the Fourier transform and the Fourier cotransform on \(L^{1}(G)\). We therefore have, for \(f \in L^{1}(G)\) and \(\chi \in \widehat{G}\):

\[
\mathcal {F} _ {\mathrm{G}} (f) (\chi) = \int_ {\mathrm{G}} \overline {{\langle \chi , x \rangle}} f (x) d x, \quad \overline {{\mathcal {F}}} _ {\mathrm{G}} (f) (\chi) = \int_ {\mathrm{G}} \langle \chi , x \rangle f (x) d x.\tag{14}
\]

In particular, \(\mathcal{F}_{\mathrm{G}}(f) = \overline{\overline{\mathcal{F}}_{\mathrm{G}}(\overline{f})}\). We also have

\[
\overline {{\mathcal {F}}} (f) (\chi) = \chi (f)\tag{15}
\]

for all \(f \in \mathrm{L}^{1}(\mathrm{G})\) and every \(\chi \in \widehat{G}\).

Let \(\sigma\) be an automorphism of G and \(\Delta\) the modulus of \(\sigma\) (INT, VII, §1, no. 4, def. 4). For \(f \in \mathrm{L}^1(\mathrm{G})\), one has

\[
\mathcal {F} (f \circ \sigma) = \Delta^ {- 1} \mathcal {F} (f) \circ \widehat {\sigma} ^ {- 1}\tag{16}
\]

(cf. loc. cit., formula (31)).

If one identifies \(\widehat{G}\) and \(X(L^{1}(G))\) (prop. 1 of II, p.~202), the Fourier cotransformation is none other than the Gelfand transformation of the Banach algebra \(L^{1}(G)\) (I, p.~7, def. 5).

PROPOSITION 4. --- The Fourier transform and the Fourier cotransform are injective morphisms of involutive algebras from L \(^{1}\) (G) into the algebra \(\mathcal{C}_{0}(\widehat{\mathrm{G}})\) of continuous functions vanishing at infinity on \(\widehat{\mathrm{G}}\).

The Fourier cotransform is a morphism of involutive algebras from \(L^{1}(G)\) into the algebra of bounded continuous functions on \(\widehat{G}\). Since it is identified with the Gelfand transform, its image is contained in \(\mathcal{C}_{0}(\widehat{\mathrm{G}})\) (I, p.~37, prop. 5), and its kernel is the radical of \(L^{1}(G)\) (prop. 8 of I, p.~38), which is zero (cor. of prop. 22 of I, p.~126).

We shall see later (corollary of Proposition 13 of II, p.~221) that the Fourier cotransformation on \(\mathcal{M}^{1}(\mathrm{G})\) is also injective.

Remark. --- The Fourier transform on the space \(L^{1}(G)\) depends on the choice of the Haar measure dx, unlike the Fourier transform on \(\mathcal{M}^{1}(G)\). If one replaces dx by the measure \(a \cdot dx\) (with \(a > 0\)), then for any integrable function f on G, the Fourier transform of f is \(a\widehat{f}\), where \(\widehat{f}\) is the Fourier transform defined with respect to the measure dx.

Consider the stellar algebra Stell(G) of the group G (def. 9 of I, p.~125), and identify L \(^{1}\) (G) with a dense subalgebra of Stell(G) (prop. 22 of I, p.~126).

PROPOSITION 5. --- By continuity, the Fourier transform and the Fourier cotransform extend uniquely to isomorphisms of stellar algebras from Stell(G) onto \(\mathcal{C}_{0}(\widehat{\mathrm{G}})\).

The Fourier cotransform extends by continuity to a morphism of stellar algebras from Stell(G) into \(\mathcal{C}_{0}(\widehat{\mathrm{G}})\). If one identifies \(\widehat{G}\) with \(\mathrm{X}(\mathrm{Stell}(\mathrm{G}))\) (cor. 1 of II, p.~204 and prop. 1 of II, p.~202), this extension is the Gelfand transform of Stell(G). By th. 1 of I, p.~108, it is an isomorphism. The assertion concerning the Fourier transform follows from this.

We will always denote by \(\overline{\mathcal{F}}\) and \(\mathcal{F}\) the isomorphisms in Proposition 5.

COROLLARY. --- \(L\), the image of \(\mathrm{L}^{1}(\mathrm{G})\) under the Fourier transform of G, is dense in \(\mathcal{C}_{0}(\widehat{\mathrm{G}})\).

Since \(L^{1}(G)\) is dense in \(\text{Stell}(G)\), this follows from Proposition 5.

PROPOSITION 6. — Suppose that G is compact. The normalized Haar measure dx belongs to \(\mathcal{M}^{1}(\mathrm{G})\), and its Fourier transform is \(\varphi_{e}\), the characteristic function of \(\{e\}\).

Let \(\chi \in \widehat{\mathrm{G}}\). Since \(\varepsilon_y * dx = dx\) for all \(y \in \mathrm{G}\), one has

\[
\mathcal {F} (d x) (\chi) = \overline {{\langle \chi , y \rangle}} \mathcal {F} (d x) (\chi)
\]

by formula (11). If \(\chi \neq 1\) , there exists \(y \in \mathrm{G}\) such that \(\langle \chi, y \rangle \neq 1\) , hence \(\mathcal{F}(dx)(\chi) = 0\) . If \(\chi = 1\) , then \(\mathcal{F}(dx)(\chi) = \int_{\mathrm{G}} dx = 1\) since the measure \(dx\) is normalized.

\subsection{Plancherel's theorem}

We denote by A(G) the vector subspace of \(L^{1}(G)\) generated by the functions \(f * g\) for \(f, g \in L^{1}(G) \cap L^{2}(G)\).

PROPOSITION 7. --- The space A(G) is a self-adjoint ideal of L¹(G). It is contained in L¹(G) ∩ L²(G), and in the image of C(G) in L¹(G).

Let \(f \in \mathrm{L}^1(\mathrm{G})\). For every \(g \in \mathrm{L}^2(\mathrm{G})\), one has \(f * g \in \mathrm{L}^1(\mathrm{G})\). Therefore, the space \(\mathrm{L}^1(\mathrm{G}) \cap \mathrm{L}^2(\mathrm{G})\) is an ideal of \(\mathrm{L}^1(\mathrm{G})\), and the same is true of the space \(\mathrm{A}(\mathrm{G})\). The ideal \(\mathrm{A}(\mathrm{G})\) is self-adjoint.

Let \(f\) and \(g\) be in \(\mathrm{L}^2(\mathrm{G})\). The convolution product \(f * g\) is then the class of the continuous function given by

\[
y \mapsto \int_ {\mathrm{G}} f (y x ^ {- 1}) g (x) d x
\]

The second assertion follows from it.

Since \(\chi(f*g) = (\chi f)*( \chi g)\) for \(\chi \in \widehat{\mathrm{G}}\), \(f \in \mathrm{L}^1(\mathrm{G})\), and \(g \in \mathrm{L}^1(\mathrm{G})\) (INT, VIII, §3, no. 1, prop. 6), we have \(\chi h \in \mathrm{A}(\mathrm{G})\) for all \(h \in \mathrm{A}(\mathrm{G})\) and \(\chi \in \widehat{\mathrm{G}}\). Since \(\varepsilon_x * f = \gamma(x)f\) and the linear representation \(\gamma\) is isometric on \(\mathrm{L}^p(\mathrm{G})\) for all \(p\) (INT, VIII, §2, no. 5, prop. 8), we have \(\varepsilon_x * f \in \mathrm{A}(\mathrm{G})\) for all \(x \in \mathrm{G}\) and \(f \in \mathrm{A}(\mathrm{G})\).

We denote by \(\widehat{\mathrm{A}}(\mathrm{G})\) the image of \(\mathrm{A}(\mathrm{G})\) under the Fourier transform. It is a subspace of \(\mathcal{C}_0(\widehat{\mathrm{G}})\).

PROPOSITION 8. --- There exists a filter base \(\mathfrak{B}\) on \(\mathrm{A}(\mathrm{G}) \cap \mathcal{K}_{+}(\mathrm{G})\) such that the following conditions are satisfied:

\begin{enumerate}
\def\labelenumi{(\roman{enumi})}
\item
  For every element \(\varphi\) of a member of \(\mathfrak{B}\), one has \(\|\varphi\|_1 = 1\) and \(\|\mathcal{F}(\varphi)\|_{\infty} \leqslant 1\);
\item
  We have
\end{enumerate}

\[
\lim _ {\varphi , \mathfrak {B}} \varphi \cdot d x = \varepsilon_ {e}
\]

in the space \(\mathcal{C}'(\mathrm{G})\) of measures with compact support on G, endowed with the topology of uniform convergence on compact subsets of \(\mathcal{C}(\mathrm{G})\);

\begin{enumerate}
\def\labelenumi{(\roman{enumi})}
\setcounter{enumi}{2}
\tightlist
\item
  We have
\end{enumerate}

\[
\lim _ {\varphi , \mathfrak {B}} \mathcal {F} (\varphi) = 1
\]

for the topology of compact convergence on \(\widehat{\mathbf{G}}\) ;

\begin{enumerate}
\def\labelenumi{(\roman{enumi})}
\setcounter{enumi}{3}
\tightlist
\item
  For \(p = 1\) or \(p = 2\), and for every \(f \in \mathrm{L}^p(\mathrm{G})\), one has \(\varphi * f \in \mathrm{A}(\mathrm{G})\) for every \(\varphi\) belonging to a member of \(\mathfrak{B}\), and
\end{enumerate}

\[
\lim _ {\varphi , \mathfrak {B}} \varphi * f = f
\]

in L \(^{p}\) (G).

Let \(K_{0}\) be a fixed compact neighborhood of e in G. Let \(B_{0}\) be a base of the neighborhood filter of e in G consisting of compact symmetric neighborhoods contained in \(K_{0}\) (cf.~TG, III, p.~4). For \(K \in B_{0}\), let \(X_{K}^{\prime}\) be the set of functions \(\psi \in \mathcal{K}_{+}(G)\) such that \(\operatorname{Supp}(\psi) \subset K\) and \(\int \psi(x)dx = 1\); it is nonempty (Lemma 1 of II, p.~200). Let \(X_{K}\) be the set of functions \(\psi * \psi\) for \(\psi \in X_{K}^{\prime}\). It is nonempty and contained in \(A(G) \cap \mathcal{K}_{+}(G)\). The set \(\mathfrak{B}\) whose elements are the sets \(X_{K}\) for K varying in \(B_{0}\) is a filter base on \(A(G) \cap \mathcal{K}_{+}(G)\). We prove that \(\mathfrak{B}\) satisfies the required properties.

If \(X \in \mathfrak{B}\) and \(\varphi \in X\), one has \(\| \varphi \|_1 = \int_{G} \varphi(x) dx = 1\), hence \(\| \mathcal{F}(\varphi) \|_\infty \leqslant 1\), which establishes property (i).

Property (ii) follows from INT, VIII, § 2, no. 7, corollary 1 of lemma 4. A compact subset of \(\widehat{\mathbf{G}}\) is a compact subset of \(\mathcal{C}(\mathbf{G})\), hence (ii) implies \(\lim_{\varphi, \mathfrak{B}} \mathcal{F}(\varphi) = 1\) for the topology of compact convergence on \(\widehat{\mathbf{G}}\), that is, (iii).

Finally, let \(p=1\) or \(p=2\). Let \(f\in\mathrm{L}^{p}(\mathrm{G})\). We have \(\varphi*f\to f\) in \(\mathrm{L}^{p}(\mathrm{G})\) according to the filter \(\mathfrak{B}\) (INT, VIII, §4, no. 7, prop. 20). Moreover, for every K in \(B_{0}\) and \(\varphi\in X_{K}\), there exists \(\psi\in X_{K}^{\prime}\) such that \(\varphi=\psi*\psi\), whence \(\varphi*f=\psi*(\psi*f)\). We have \(\psi\in L^{1}(G)\cap L^{2}(G)\) and \(\psi*f\in L^{1}(G)\cap L^{2}(G)\), hence \(\varphi*f\in A(G)\).

COROLLARY 1. --- The space A(G) is dense in L \(^{1}\) (G) and in L \(^{2}\) (G). It is also dense in Stell(G), and its image \(\widehat{\mathrm{A}}(G)\) under the Fourier transformation is dense in \(\mathcal{C}_{0}(\widehat{\mathrm{G}})\).

Assertion (iv) of the proposition provides the first assertion. Since L \(^{1}\) (G) is dense in Stell(G), the second follows from it, and the last then follows from prop. 5 of II, p.~209.

COROLLARY 2. --- For \(f \in \mathrm{A}(\mathrm{G})\), let \(\Omega_f\) be the set of \(\chi \in \widehat{\mathrm{G}}\) such that \(\mathcal{F}(f)(\chi) \neq 0\). The sets \(\Omega_f\) form an open covering of \(\widehat{\mathrm{G}}\).

This follows from the preceding corollary since, for every \(\chi \in \widehat{\mathrm{G}}\), the map \(f \mapsto \mathcal{F}(f)(\chi)\) is a nonzero character of \(\mathrm{L}^1(\mathrm{G})\).

Recall that the left regular representation \(\pmb{\gamma}\) of Stell(G) on \(\mathrm{L}^2 (\mathrm{G})\) (cf.~I, p.~125, no. 13) is denoted \(\pmb{\gamma}(\varphi)f = \varphi *f\) for \(\varphi \in \mathrm{Stell}(\mathrm{G})\) and \(f\in \mathrm{L}^2 (\mathrm{G})\).

Lemma 3. --- For every \(f \in \mathrm{A}(\mathrm{G})\), there exists a unique bounded measure \(\mu_{f}\) on \(\widehat{\mathrm{G}}\) such that

\[
(\varphi * f) (e) = \int_ {\widehat {G}} \mathcal {F} (\varphi) d \mu_ {f}\tag{17}
\]

for all \(\varphi\in\mathrm{Stell}(\mathrm{G})\).

Moreover, for all \(f\) and \(g\) in A(G), we have the equality

\[
\mathcal {F} (f) \cdot \mu_ {g} = \mathcal {F} (g) \cdot \mu_ {f},\tag{18}
\]

between the measure with density \(\mathcal{F}(f)\) with respect to \(\mu_g\) and the measure with density \(\mathcal{F}(g)\) with respect to \(\mu_f\).

Let \(f\) and \(g\) be elements of \(\mathrm{L}^{1}(\mathrm{G}) \cap \mathrm{L}^{2}(\mathrm{G})\). For every \(\varphi \in \mathrm{Stell}(\mathrm{G})\), one has \(\varphi * f \in \mathrm{L}^{2}(\mathrm{G})\) and \(\| \varphi * f \|_{2} \leqslant \| \varphi \|_{*} \| f \|_{2}\) (I, p.~126, formula (8)). Moreover, we have \(\varphi * (f * g) = (\varphi * f) * g\) (loc. cit., formula (9)). This latter function belongs to the closure \(\mathcal{C}_{0}(\mathrm{G})\) of \(\mathcal{K}(\mathrm{G})\) in \(\mathcal{C}(\mathrm{G})\) (INT, VIII, §4, n° 5, prop. 15). Moreover, we have

\[
\left\| \varphi * (f * g) \right\| _ {\infty} \leqslant \left\| \varphi * f \right\| _ {2} \| g \| _ {2} \leqslant \left\| \varphi \right\| _ {*} \| f \| _ {2} \| g \| _ {2}.
\]

Since the functions \(f * g\) for \(f\) and \(g\) in \(\mathrm{L}^1(\mathrm{G}) \cap \mathrm{L}^2(\mathrm{G})\) generate \(\mathrm{A}(\mathrm{G})\), one deduces that \(\varphi * f \in \mathcal{C}_0(\mathrm{G})\) for all \(f \in \mathrm{A}(\mathrm{G})\) and \(\varphi \in \mathrm{Stell}(\mathrm{G})\), and that the map \(\varphi \mapsto (\varphi * f)(e)\) is a continuous linear form on Stell(G). Since \(\mathcal{F}\) is an isomorphism from Stell(G) onto \(\mathcal{C}_{0}(\widehat{\mathrm{G}})\) (prop. 5 of II, p.~209), the first assertion follows from it.

Now let \(f\) and \(g\) be in \(\mathrm{A}(\mathrm{G})\). For \(\varphi \in \mathrm{L}^1 (\mathrm{G})\), one has

\[
\begin{array}{r l r} & & {\big (\mathcal {F} (f) \cdot \mu_ {g} \big) (\mathcal {F} (\varphi)) = \int_ {\widehat {\mathrm{G}}} \mathcal {F} (\varphi) \mathcal {F} (f) d \mu_ {g} = \int_ {\widehat {\mathrm{G}}} \mathcal {F} (\varphi * f) d \mu_ {g}} \\ & & {= ((\varphi * f) * g) (e).} \end{array}\tag{19}
\]

Since \((\varphi*f)*g=(\varphi*g)*f\) and since the image of \(L^{1}(G)\) under the Fourier transform is dense in \(\mathcal{C}_{0}(\widehat{\mathrm{G}})\) (cor. of II, p.~210), one deduces from formula (19) that formula (18) is satisfied for all \(f\) and \(g\) in \(A(G)\).

Lemma 4. --- There exists a unique measure \(\nu\) on \(\widehat{\mathbf{G}}\) such that

\[
\mu_ {f} = \mathcal {F} (f) \cdot \nu
\]

for all \(f \in \mathrm{A}(\mathrm{G})\). For \(f \in \mathrm{A}(\mathrm{G})\), one has \(\mathcal{F}(f) \in \mathrm{L}^{1}(\widehat{\mathrm{G}}, \nu) \cap \mathrm{L}^{2}(\widehat{\mathrm{G}}, \nu)\). Let \(f \in \mathrm{A}(\mathrm{G})\). Denote by \(\Omega_{f}\) the open set in \(\widehat{G}\) consisting of the \(\chi \in \widehat{G}\) such that \(\mathcal{F}(f)(\chi) \neq 0\). Let \(\varphi\) be the characteristic function of \(\widehat{G} - \Omega_{f}\). For every \(g \in \mathrm{A}(\mathrm{G})\), we then have

\[
\int_ {\widehat {\mathrm{G}}} \mathcal {F} (g) d (\boldsymbol {\varphi} \cdot \boldsymbol {\mu} _ {f}) = \int_ {\widehat {\mathrm{G}}} \varphi \mathcal {F} (f) d \mu_ {g} = 0
\]

given formula (18).

By Corollary 1 of II, p.~212, the image \(\widehat{\mathrm{A}}(\mathrm{G})\) of \(\mathrm{A}(\mathrm{G})\) under the Fourier transform is dense in \(\mathcal{C}_0(\widehat{\mathrm{G}})\). We then deduce from the preceding formula that \(\varphi \cdot \mu_f = 0\), so that \(\mu_f\) is concentrated on \(\Omega_f\) (INT, IV, §4, no. 7, def. 4). Let \(\nu_f\) be the measure on \(\Omega_f\) with density \(\mathcal{F}(f)^{-1}\) with respect to \(\mu_f|\Omega_f\).

The sets \(\Omega_f\), for \(f \in \mathrm{A}(\mathrm{G})\), form an open covering of \(\widehat{\mathrm{G}}\) (cor. 2). For all \(f\) and \(g\) in \(\mathrm{A}(\mathrm{G})\), formula (18) shows that \(\nu_f|(\Omega_f \cap \Omega_g) = \nu_g|(\Omega_f \cap \Omega_g)\). Consequently, there exists a unique measure \(\nu\) on \(\widehat{\mathrm{G}}\) such that \(\nu_f = \nu|\Omega_f\) for all \(f \in \mathrm{A}(\mathrm{G})\) (INT, III, §2, no. 1, prop. 1).

If \(f \in \mathrm{A}(\mathrm{G})\), the measures \(\mu_f\) and \(\mathcal{F}(f) \cdot \nu\) are concentrated on \(\Omega_f\), and their restrictions to \(\Omega_f\) are equal to \(\mathcal{F}(f) \cdot \nu_f\); these measures are therefore equal.

Since \(\mu_f\) is a bounded measure, the Fourier transform \(\mathcal{F}(f)\) belongs to the space \(\mathrm{L}^1 (\widehat{\mathrm{G}},\nu)\). Moreover, since \(\mathcal{F}(f)\) belongs to \(\mathcal{C}_0(\widehat{\mathrm{G}})\), we also have \(\mathcal{F}(f)\in \mathrm{L}^2 (\widehat{\mathrm{G}},\nu)\).

Formula (17) now reads, for \(\varphi\in\mathrm{Stell}(\mathrm{G})\) and \(f\in\mathrm{A}(\mathrm{G})\):

\[
(\varphi * f) (e) = \int_ {\widehat {G}} \mathcal {F} (\varphi) \mathcal {F} (f) d \nu .\tag{20}
\]

In particular, for \(f\) and \(g\) in \(\mathrm{A}(\mathrm{G})\), one has

\[
\int_ {\widehat {\mathrm{G}}} \mathcal {F} (f) \mathcal {F} (g) d \nu = (f * g) (e) = \int_ {\mathrm{G}} f (x) g (x ^ {- 1}) d x.\tag{21}
\]

PROPOSITION 9. — The measure \(\nu\) characterized by Lemma 4 is a Haar measure on \(\widehat{\mathbf{G}}\).

Let \(\chi \in \widehat{\mathrm{G}}\). For \(f\) and \(g\) in \(\mathrm{A}(\mathrm{G})\), let us apply formula (21) to \(\chi f\) and \(\chi g\). The right-hand side is unchanged, hence

\[
\nu (\mathcal {F} (f) \mathcal {F} (g)) = \nu (\mathcal {F} (\chi f) \mathcal {F} (\chi g)).
\]

From formulas (12) of II, p.~208 and (13) of II, p.~208, it follows that

\[
\nu (\mathcal {F} (f) \mathcal {F} (g)) = \nu (\varepsilon_ {\chi} * (\mathcal {F} (f) \mathcal {F} (g))).
\]

We then deduce from INT, VIII, §4, n° 3, prop. 7 that

\[
\nu (\mathcal {F} (f) \mathcal {F} (g)) = (\varepsilon_ {\chi^ {- 1}} * \nu) (\mathcal {F} (f) \mathcal {F} (g)),
\]

that is,

\[
(\mathscr {F} (f) \cdot \nu) (\mathscr {F} (g)) = (\mathscr {F} (f) \cdot (\varepsilon_ {\chi^ {- 1}} * \nu)) (\mathscr {F} (g)).
\]

As the space \(\widehat{\mathrm{A}}(\mathrm{G})\) is dense in \(\mathcal{C}_{0}(\widehat{\mathrm{G}})\) (cor. 1), it follows that equality holds

\[
\mathscr {F} (f) \cdot \nu = \mathscr {F} (f) \cdot (\varepsilon_ {\chi^ {- 1}} * \nu).
\]

The measures \(\nu\) and \(\varepsilon_{\chi^{-1}} * \nu\) therefore coincide on the open set \(\Omega_f\) where \(\mathcal{F}(f)\) is nonzero. By Corollary 2 and INT, III, §2, no. 1, cor. of Prop. 1, these measures are therefore equal.

This shows that the measure \(\nu\) is proportional to a Haar measure on \(\widehat{\mathrm{G}}\). Let \(f \in \mathrm{A}(\mathrm{G})\). Let us take \(g = \widetilde{f}\) in formula (21). It is then written as

\[
\int_ {\widehat {\mathrm{G}}} | \mathcal {F} (f) | ^ {2} d \nu = \int_ {\mathrm{G}} | f | ^ {2} d x,\tag{22}
\]

which proves that the measure \(\nu\) is not zero. The measure \(\nu\) is therefore a Haar measure on \(\widehat{G}\).

DEFINITION 4. --- The Haar measure \(\nu\) on \(\widehat{G}\) from Proposition 9 is called the dual measure of the given Haar measure dx on G.

We will often denote by \(d\chi\) or \(d\widehat{x}\) the Haar measure on \(\widehat{G}\) which is dual to the Haar measure dx.

Note. --- Let \(a\) be a real number \(>0\). If one replaces \(dx\) by the measure \(a \cdot dx\), the convolution product of the functions \(f\) and \(g \in \mathrm{L}^{1}(\mathrm{G})\) is replaced by \(a(f*g)\). We have seen (II, p.~209, remark) that \(\mathcal{F}(f)\) is replaced by \(a\mathcal{F}(f)\). Therefore \(\mu_{f}\) is unchanged and \(\nu\) is replaced by \(a^{-1} \cdot \nu\). In particular, the measure \(dx \otimes d\hat{x}\) on \(\mathrm{G} \times \hat{\mathrm{G}}\) is independent of the choice of Haar measure on \(\mathrm{G}\).

Lemma 5. --- The space A(\(\hat{G}\)) is dense in L²(\(\hat{G}\)).

Let \(h\) be an element of \(\mathrm{L}^2(\widehat{\mathrm{G}})\) orthogonal to \(\widehat{\mathrm{A}}(\mathrm{G})\). For \(f\) and \(g\) in \(\mathrm{A}(\mathrm{G})\), one has \(\mathcal{F}(f) \cdot \mathcal{F}(g) = \mathcal{F}(f*g) \in \widehat{\mathrm{A}}(\mathrm{G})\), hence \(h \cdot \overline{\mathcal{F}(f)}\) is orthogonal to \(\mathcal{F}(g)\). Thus, for every \(f \in \mathrm{A}(\mathrm{G})\), the function \(h \cdot \mathcal{F}(f)\) is orthogonal to \(\widehat{\mathrm{A}}(\mathrm{G})\). But \(h \cdot \mathcal{F}(f) \in \mathrm{L}^1(\widehat{\mathrm{G}})\), and \(\widehat{\mathrm{A}}(\mathrm{G})\) is dense in \(\mathcal{C}_0(\widehat{\mathrm{G}})\), so the measure \(h\mathcal{F}(f) \cdot \nu\) is zero, that is, \(h\mathcal{F}(f)\) is \(\nu\)-locally negligible (INT, V, §5, no. 3, cor. 2 of prop. 3). In particular, \(h\) is \(\nu\)-locally negligible on the set \(\Omega_f\) of characters \(\chi\) such that \(\mathcal{F}(f)(\chi) \neq 0\). By corollary 2, it follows that \(h\) is \(\nu\)-locally negligible, hence zero since \(h\) belongs to \(\mathrm{L}^2(\widehat{\mathrm{G}})\). This concludes the proof.

THEOREM 1 (Plancherel). --- The restriction of the Fourier transform to the subspace A(G) of L²(G) extends uniquely to an isometry Φ of L²(G) onto L²(Ĝ).

Moreover, if \(f \in \mathrm{L}^{1}(\mathrm{G}) \cap \mathrm{L}^{2}(\mathrm{G})\), its Fourier transform belongs to \(\mathrm{L}^{2}(\widehat{\mathrm{G}})\) and coincides in \(\mathrm{L}^{2}(\widehat{\mathrm{G}})\) with \(\Phi(f)\).

By formula (22), the restriction of \(\mathcal{F}\) to A(G) is an isometry of the subspace A(G) of L \(^2\) (G) onto the subspace \(\hat{\mathrm{A}}(\mathrm{G})\) of L \(^2\) (\(\hat{\mathrm{G}}\)). Since A(G) is dense in L \(^2\) (G) (cor. 1 of II, p.~212), the Fourier transform extends uniquely to an isometry \(\Phi\) of L \(^2\) (G) onto a closed subspace of L \(^2\) (\(\hat{\mathrm{G}}\)). But since its image contains \(\hat{\mathrm{A}}(\mathrm{G})\), which is dense in L \(^2\) (\(\hat{\mathrm{G}}\)) (lemma 5), the map \(\Phi\) is surjective.

Now let \(f \in \mathrm{L}^{1}(\mathrm{G}) \cap \mathrm{L}^{2}(\mathrm{G})\); let us show that its Fourier transform belongs to \(\mathrm{L}^{2}(\widehat{\mathrm{G}})\). By prop. 8, (iv) of II, p.~211, and the fact that \(\mathrm{A}(\mathrm{G})\) is an ideal of \(\mathrm{L}^{1}(\mathrm{G})\), there exists a filter base B on \(\mathrm{A}(\mathrm{G})\) which converges to f both in \(\mathrm{L}^{1}(\mathrm{G})\) and in \(\mathrm{L}^{2}(\mathrm{G})\). We then have

\[
\Phi (f) = \lim _ {g, \mathfrak {B}} \Phi (g) = \lim _ {g, \mathfrak {B}} \mathcal {F} (g)
\]

in \(\mathrm{L}^{2}(\widehat{\mathrm{G}})\), and \(\mathcal{F}(f)=\lim_{g,\mathfrak{B}}\mathcal{F}(g)\) in \(\mathcal{C}_{0}(\widehat{\mathrm{G}})\). There therefore exists a sequence \((g_{n})\) in \(\mathrm{A}(\mathrm{G})\) such that \(\mathcal{F}(g_{n})\) converges to \(\Phi(f)\) in \(\mathrm{L}^{2}(\widehat{\mathrm{G}})\) and to \(\mathcal{F}(f)\) in \(\mathcal{C}_{0}(\widehat{\mathrm{G}})\). By INT, IV, §3, no. 4, th. 3 and cor. 1, we have \(\mathcal{F}(f)=\Phi(f)\), and in particular \(\mathcal{F}(f)\in\mathrm{L}^{2}(\widehat{\mathrm{G}})\). This completes the proof of the theorem.

We still denote by \(\mathcal{F}\) the isometry of \(\mathrm{L}^2 (\mathrm{G})\) onto \(\mathrm{L}^2 (\widehat{\mathrm{G}})\) defined in Theorem 1, and we call it the Fourier transformation on \(\mathrm{L}^2 (\mathrm{G})\). Likewise, the Fourier cotransform admits a unique isometric extension to \(\mathrm{L}^2 (\mathrm{G})\), still called the Fourier cotransform and denoted \(\overline{\mathcal{F}}\).

COROLLARY. --- Suppose that G is compact and that dx is the normalized Haar measure on G. Then the family of unitary characters of G is an orthonormal basis of L²(G).

Since G is compact, the characters of G belong to \(L^{2}(G)\). For \(\chi\) and \(\xi\) in \(\widehat{G}\), one has

\[
\int_ {\mathrm{G}} \overline {{\langle \chi , x \rangle}} \langle \xi , x \rangle d x = \mathcal {F} _ {\mathrm{G}} (d x) (\chi \xi^ {- 1}),
\]

so the family of characters of G is orthonormal (prop. 6 of II, p.~210). It is moreover total because the scalar product of \(\chi \in \widehat{\mathbf{G}}\) and of \(f \in \mathrm{L}^2(\mathrm{G})\) is equal to \(\mathcal{F}_{\mathrm{G}}(f)(\chi)\), and hence \(f\) is orthogonal to \(\widehat{\mathbf{G}}\) if and only if \(\mathcal{F}_{\mathrm{G}}(f)\) is zero in \(\mathrm{L}^2(\widehat{\mathbf{G}})\), if and only if \(f\) is zero (th. 1).

Remarks. --- 1) Some of the formulas concerning the Fourier transformation on \(\mathrm{L}^{1}(\mathrm{G})\) extend to the Fourier transformation on \(\mathrm{L}^{2}(\mathrm{G})\). In particular, for \(f \in \mathrm{L}^{2}(\mathrm{G})\) and \(\chi \in \widehat{\mathbf{G}}\), one has

\[
\overline {{\mathcal {F}}} (f) = (\chi \mapsto \mathcal {F} (f) (\chi^ {- 1})) = \overline {{\mathcal {F} (\overline {{f}})}}
\]

\[
\mathcal {F} (\varepsilon_ {x} * f) = \eta (x ^ {- 1}) \mathcal {F} (f), \qquad \overline {{\mathcal {F}}} (\varepsilon_ {x} * f) = \eta (x) \mathcal {F} (f)
\]

\[
\mathcal {F} (\chi f) = \varepsilon_ {\chi} * \mathcal {F} (f), \quad \overline {{\mathcal {F}}} (\chi f) = \varepsilon_ {\chi} * \mathcal {F} (f).
\]

If \(\sigma\) is an automorphism of G and \(\Delta\) the modulus of \(\sigma\) (INT, VII, §1, no. 4, def. 4), then for \(f \in \mathrm{L}^2(\mathrm{G})\), one has

\[
\mathcal {F} (f \circ \sigma) = \Delta^ {- 1} \mathcal {F} (f) \circ \widehat {\sigma} ^ {- 1}
\]

in L \(^{2}\) (G).

\begin{enumerate}
\def\labelenumi{\arabic{enumi})}
\setcounter{enumi}{1}
\tightlist
\item
  The formulas
\end{enumerate}

\[
\| \overline {{\mathcal {F}}} (f) \| ^ {2} = \| f \| ^ {2}\tag{23}
\]

for \(f\in \mathrm{L}^2 (\mathrm{G})\), or

\[
\int_ {\mathrm{G}} f (x) \overline {{g (x)}} d x = \int_ {\widehat {\mathrm{G}}} \mathcal {F} (f) (\chi) \overline {{\mathcal {F} (g) (\chi)}} d \chi\tag{24}
\]

for \(f\) and \(g\) in \(\mathrm{L}^2 (\mathrm{G})\), are called “Plancherel formulas”.

PROPOSITION 10. --- Let \(n \geqslant 0\) be an integer and let \(\mathrm{G}_{1}, \cdots, \mathrm{G}_{n}\) be locally compact abelian groups. Let \(\mu_{j}\), for \(1 \leqslant j \leqslant n\), be Haar measures on \(\mathrm{G}_{j}\). Let G be the product group of the groups \(\mathrm{G}_{j}\) for \(1 \leqslant j \leqslant n\). Let \(\beta\) be the isomorphism from \(\widehat{\mathrm{G}}\) onto \(\prod \widehat{\mathrm{G}}_{j}\) of prop. 2 of II, p.~206. The Haar measure on \(\widehat{\mathrm{G}}\) dual to the product Haar measure \(\mu = \mu_{1} \otimes \cdots \otimes \mu_{n}\) on G is identified with the product of the Haar measures \(\widehat{\mu}_{j}\).

Indeed, for every family \((f_{j})\) of nonzero elements of \(\mathcal{L}^{2}(\mathrm{G}_{j})\), the function f on G defined by \((x_{j})\mapsto\Pi f_{j}(x_{j})\) belongs to \(\mathcal{L}^{2}(\mathrm{G})\), and satisfies

\[
\int_ {\mathrm{G}} | f | ^ {2} d \mu = \prod_ {j} \int_ {\mathrm{G} _ {j}} | f _ {j} | ^ {2} d \mu_ {j} = \prod_ {j} \int_ {\widehat {\mathrm{G}} _ {j}} | \mathcal {F} _ {\mathrm{G} _ {j}} (f _ {j}) | ^ {2} d \widehat {\mu} _ {j},
\]

by the Plancherel formula, which proves that the product Haar measure of the \(\widehat{\mu}_{j}\) is identified with the dual Haar measure of \(\mu\).

\subsection{The Fourier inversion formula}

Recall that every element \(f\) of A(G) is the class of a unique continuous function (prop. 7, b) of II, p.~210). For \(x \in \mathrm{G}\), we shall denote by \(f(x)\) the value of this function at \(x\).

PROPOSITION 11. --- Let \(f \in \mathrm{A}(\mathrm{G})\). Then \(\mathcal{F}(f) \in \mathrm{L}^1(\widehat{\mathrm{G}})\) and, for every \(x \in \mathrm{G}\), one has

\[
f (x) = \int_ {\widehat {G}} \langle \widehat {x}, x \rangle \mathcal {F} (f) (\widehat {x}) d \widehat {x}.\tag{25}
\]

In other words, for \(f \in \mathrm{A}(\mathrm{G})\), one has

\[
f = \overline {{\mathcal {F}}} _ {\widehat {\mathrm{G}}} (\mathcal {F} _ {\mathrm{G}} (f)) \circ \eta ,\tag{26}
\]

where \(\eta\) denotes the canonical mapping from G into the bidual group \(\widehat{\widehat{G}}\).

By lemma 4 of II, p.~213 and proposition 9 of II, p.~214, we have \(\mathcal{F}(f) \in \mathrm{L}^1(\widehat{\mathrm{G}})\) for every function \(f \in \mathrm{A}(\mathrm{G})\). By the Plancherel formula (24), for \(f\) and \(g\) in \(\mathrm{L}^2(\mathrm{G})\), one has

\[
(f * \widetilde {g}) (e) = \int_ {\widehat {\mathrm{G}}} \mathcal {F} (f) (\chi) \overline {{\mathcal {F} (g) (\chi)}} d \widehat {x} (\chi).\tag{27}
\]

Let \(f\) and \(g\) be in \(\mathrm{L}^1(\mathrm{G}) \cap \mathrm{L}^2(\mathrm{G})\), and let \(h = f * \widetilde{g} \in \mathrm{A}(\mathrm{G})\). Since the Fourier transformation is an involutive morphism, formula (27) is assertion (25) for the function \(h\) at the point \(x = e\). By linearity, we deduce that formula (25) is valid at the point \(x = e\) for every function \(h \in \mathrm{A}(\mathrm{G})\).

Let \(x \in G\) and \(h \in A(G)\). Let \(h_{1} = \varepsilon_{x^{-1}} * h\). Then \(h_{1} \in A(G)\) and \(h_{1}(e) = h(x)\). Since moreover \(\mathcal{F}(h_{1})(\chi) = \overline{\langle\chi, x\rangle}\mathcal{F}(f)(\chi)\) for all \(\chi \in \widehat{G}\) (cf. formula (11) of II, p. 208), formula (25) for the function \(h_{1}\) at the point e implies formula (25) for h at the point x.

Lemma 6. — Let \(\varphi\in\mathrm{L}^{1}(\widehat{\mathrm{G}})\cap\mathrm{L}^{2}(\widehat{\mathrm{G}})\). Then \(f=\overline{\mathcal{F}}_{\widehat{\mathrm{G}}}(\varphi)\circ\eta\) belongs to \(\mathrm{L}^{2}(\mathrm{G})\), and \(\mathcal{F}_{\mathrm{G}}(f)=\varphi\) in \(\mathrm{L}^{2}(\widehat{\mathrm{G}})\).

The function \(f\) is continuous and bounded on \(G\) because \(\varphi \in L^{1}(\widehat{G})\). For every function \(g \in L^{1}(G) \cap L^{2}(G)\), one has

\[
\begin{array}{r} \int_ {\mathrm{G}} g (x) f (x) d x = \int_ {\mathrm{G}} g (x) \Big (\int_ {\widehat {\mathrm{G}}} \langle \chi , x \rangle \varphi (\chi) d \widehat {x} (\chi) \Big) d x \\ = \int_ {\widehat {\mathrm{G}}} \overline {{\mathcal {F}}} _ {\mathrm{G}} (g) (\chi) \varphi (\chi) d \widehat {x} (\chi), \end{array}\tag{28}
\]

by applying the Lebesgue-Fubini theorem (INT, V, §8, no. 4, th. 1, a)) to the function \((x,\chi)\mapsto g(x)\varphi(\chi)\langle\chi,x\rangle\), which is integrable on \(G\times\widehat{G}\) with respect to the product measure \(dx\otimes d\widehat{x}\). From this we deduce that

\[
\left| \int_ {\mathrm{G}} g (x) f (x) d x \right| \leqslant \| \overline {{{\mathcal {F}}}} _ {\mathrm{G}} (g) \| _ {2} \| \varphi \| _ {2} = \| g \| _ {2} \| \varphi \| _ {2},
\]

by the Plancherel formula. The linear form \(g \mapsto \int_{G} fg\) is therefore continuous on \(\mathrm{L}^{1}(\mathrm{G}) \cap \mathrm{L}^{2}(\mathrm{G})\), and since \(\mathrm{L}^{1}(\mathrm{G}) \cap \mathrm{L}^{2}(\mathrm{G})\) is dense in the Hilbert space \(\mathrm{L}^{2}(\mathrm{G})\), we deduce that f belongs to \(\mathrm{L}^{2}(\mathrm{G})\).

Applying th. 1 of II, p. 215, we obtain on the other hand

\[
\begin{array}{r} \int_ {\mathrm{G}} g (x) f (x) d x = \int_ {\widehat {\mathrm{G}}} \overline {{\mathcal {F} _ {\mathrm{G}} (\overline {{g}}) (\chi)}} \mathcal {F} _ {\mathrm{G}} (f) (\chi) d \widehat {x} (\chi) \\ = \int_ {\widehat {\mathrm{G}}} \overline {{\mathcal {F}}} _ {\mathrm{G}} (g) (\chi) \mathcal {F} _ {\mathrm{G}} (f) (\chi) d \widehat {x} (\chi) \end{array}
\]

for all \(g \in \mathrm{L}^{2}(\mathrm{G})\). Comparing with (28), we conclude that \(\varphi = \mathcal{F}_{\mathrm{G}}(f)\) in \(\mathrm{L}^{2}(\widehat{\mathrm{G}})\), since \(\mathrm{A}(\mathrm{G})\) is contained in \(\mathrm{L}^{1}(\mathrm{G}) \cap \mathrm{L}^{2}(\mathrm{G})\) and \(\widehat{\mathrm{A}}(\mathrm{G})\) is dense in \(\mathrm{L}^{2}(\widehat{\mathrm{G}})\) (lemma 5 of II, p.~215).

PROPOSITION 12 (Fourier inversion formula)

Let \(f \in \mathrm{L}^2(\mathrm{G})\) be such that \(\mathcal{F}_{\mathrm{G}}(f) \in \mathrm{L}^1(\widehat{\mathrm{G}})\). Then we have \(f = \overline{\mathcal{F}}_{\widehat{\mathrm{G}}}(\mathcal{F}_{\mathrm{G}}(f)) \circ \eta\) in \(\mathrm{L}^2(\mathrm{G})\). In other words, for almost all \(x \in \mathrm{G}\), we have

\[
f (x) = \int_ {\widehat {G}} \langle \widehat {x}, x \rangle \mathcal {F} _ {G} (f) (\widehat {x}) d \widehat {x}.
\]

The function \(\varphi = \mathcal{F}_{\mathrm{G}}(f)\) belongs to \(\mathrm{L}^1 (\widehat{\mathrm{G}})\cap \mathrm{L}^2 (\widehat{\mathrm{G}})\), and we obtain the desired formula by applying the lemma.

COROLLARY 1. --- For every closed subset P of \(\widehat{\mathbf{G}}\) and every \(\chi \in \widehat{\mathbf{G}}\) not belonging to P, there exists a function \(f \in \mathrm{L}^{1}(\mathrm{G})\) such that \(\mathcal{F}(f)\) is zero on P and nonzero at \(\chi\).

Since, by (12), we have \(\mathcal{F}(\chi f) = \varepsilon_{\chi} * \mathcal{F}(f)\) for all \(\chi \in \widehat{\mathrm{G}}\), it suffices to consider the case where \(\chi\) is the identity element of \(\widehat{\mathrm{G}}\).

Let U be a compact symmetric neighborhood of \(e \in \widehat{G}\) such that \(U^{2} \cap P = \varnothing\). Let \(\varphi\) be a continuous positive function on \(\widehat{G}\), vanishing outside U and such that \(\varphi(e) = 1\). The function \(\varphi_{1} = \varphi * \varphi\) is then zero on P and \(\varphi_{1}(e) > 0\). It therefore suffices to show that \(\varphi_{1}\) belongs to the image of the Fourier transform on \(L^{1}(G)\). Moreover, \(\varphi\) and \(\varphi_{1}\) belong to \(L^{1}(\widehat{G}) \cap L^{2}(\widehat{G})\). Set \(f = \overline{\mathcal{F}}(\varphi) \circ \eta\) and \(f_{1} = \overline{\mathcal{F}}(\varphi_{1}) \circ \eta\). Lemma 6 implies that f and \(f_{1}\) belong to \(L^{2}(G)\) and satisfy \(\varphi = \mathcal{F}(f)\) and \(\varphi_{1} = \mathcal{F}(f_{1})\). Moreover

\[
f _ {1} = \overline {{\mathcal {F}}} (\varphi * \varphi) \circ \eta = (\overline {{\mathcal {F}}} (\varphi) \circ \eta) ^ {2} = f ^ {2},
\]

and therefore \(f_{1} \in \mathrm{L}^{1}(\mathrm{G})\). Thus \(\varphi_{1} = \mathcal{F}(f_{1})\) is indeed in the image of \(\mathrm{L}^{1}(\mathrm{G})\) under the Fourier transform.

COROLLARY 2. --- The Banach algebra L \(^{1}\) (G) is regular (I, p.~89, def. 1).

By prop. 1 of I, p.~88 and the identification of the Gelfand transform of L \(^{1}\) (G) with the Fourier cotransform of G, this follows from the preceding corollary.
